%% 第31回編集委員会（物理第5巻）：組 原W4
%% p5-08.tex  原-8　固体の中の電子：エネルギーバンド（★2〜5・3段）
%% 中身は計画書 tools/phys5-plan.md の B「原-8」。公式の持ち主は A（A9・A21・A54・A55・A57・A84・A85）。
%%   ほかの組の章は、計画書 B の節の番号で指した（原-2 ⑨⑩、原-3 ②、原-5 ③④⑦、原-6 ⑧⑪、原-7 ④⑤⑦⑧⑪、原-9 ③⑪、原-10）。
%%   波動関数は計画書 C-1 のとおり $\psi$。虚数単位は物理の慣用の $i$（N-3）。
%% 数値はすべて Python で計算した：checks/物理第5巻/原-8/calc_p5_08.py（出力 calc_p5_08_out.txt）、
%%   code_p5_08.py（⑪の囲みのプログラム）。デバイの曲線の点は debye.py（同じフォルダ）。常用漢字は 原-7 のフォルダの kanji_p5_07_08.py。
%% 前の巻・数学の本との約束（計画書 A-6・C-6-2）
%%   ・Q45（電-4「自由電子は1秒に千 km 以上で動き回る（原-8）」）→ ④・例題4（v_F＝1.57×10⁶ m/s＝1570 km/s）
%%   ・Q48（電-4 ドルーデの τ＝2.4×10⁻¹⁴ s、「電子の動きは量子力学で（原-8）」）→ ⑩・例題9（v_F τ＝38 nm）
%%   ・Q49（電-4「抵抗率が 10²⁴ 倍違う理由はバンド（原-8・9）」）→ ⑧⑨（はてな？）
%%   ・Q16（熱-15 ⑨ フェルミ分布を原-8 で使う。金属の電子の比熱）→ ⑤⑥（熱-15 の ε・μ を E・E_F と書くと一言）
%%   ・Q14（熱-15 の seikaku「デバイの理論は原-8 で」）→ ⑪・例題11（零点エネルギーは原-6 ⑧）
%%   ・Q43（電-3 強誘電体のしくみ（原-8））→ ⑫
%%   ・Q47（電-4 熱電対の電圧を電子から（原-8・9））→ ⑫（金属の分。半導体とペルチェ素子は原-9 ⑪）
%%   ・Q53（電-6・電-8 超伝導のしくみ（原-8）、MRI）→ ⑫・例題12
%%   ・Q54（電-7 磁区・交換相互作用（原-7・8））→ ⑫（帯の中の強磁性。交換相互作用のもとは原-7 ⑧）
%%   ・電-4 の銅の n＝8.49×10²⁸ /m³・τ＝2.4×10⁻¹⁴ s・抵抗率 1.72×10⁻⁸ Ω·m（計画書 C-4 の2）→ 例題4・9・10 で同じ数
%%   ・熱-11 の銅の熱伝導率 約400 W/(m·K)→ 例題10 で比べた
%%   ・ケイ素の禁制帯 1.12 eV（計画書 C-4 の2：原-2・8・9 で同じ数）、箱 1 nm の 0.376 eV（原-5）→ ④で比べた
%% 記号（この章）：格子定数 a、面の間かく d、X線の波長 λ（第4部の崩壊定数とは部が別）、ブラッグの角 θ、
%%   電子の数密度 n（第2部の主量子数とは部が別）、フェルミエネルギー E_F（熱-15 の化学ポテンシャル μ と同じ物）、
%%   フェルミ温度 T_F、フェルミ速度 v_F、波数 k（ボルツマン定数はかならず k_B）、禁制帯幅 E_g、平均時間 τ（電-4 と同じ）、
%%   平均自由行程 ℓ、導電率 σ・熱伝導率 κ（熱-11 は k と書いたが、この章では波数とぶつかるので κ）、デバイ温度 θ_D、
%%   超伝導のすき間 Δ、くしの強さ P（電力・圧力はこの章に出ない）、結晶の中の波の番号 K（ブロッホの波数）。
%% 出典（11月11日に確かめた）
%%   ・定数（h・ħ・e・c・k_B・N_A・m_e・R）とケイ素の格子定数 a＝0.5431020511 nm：CODATA 2022
%%       https://physics.nist.gov/cuu/Constants/Table/allascii.txt（計画書 D4）
%%   ・銅の Kα1 線 8047.8227 eV（波長 0.154059 nm）：NIST X-ray Transition Energies
%%       https://physics.nist.gov/cgi-bin/XrayTrans/search.pl?element=Cu&trans=KL3（curl で読んだ）
%%   ・密度（ケイ素 2.330、銅 8.960、銀 10.50、金 19.32、アルミニウム 2.699、ナトリウム 0.971 g/cm³）：NIST X-Ray Mass Attenuation Coefficients 表1
%%       https://physics.nist.gov/PhysRefData/XrayMassCoef/tab1.html（計画書 D7）
%%   ・窒素の沸点 77.3 K：NIST Chemistry WebBook https://webbook.nist.gov/cgi/cbook.cgi?ID=C7727379&Mask=4（curl で読んだ）
%%   ・ノーベル賞：1914 ラウエ（結晶による X線回折）、1915 ブラッグ父子（X線による結晶構造の解析）、1972 バーディーン・クーパー・シュリーファー（BCS 理論）、
%%       1987 ベドノルツ・ミュラー（セラミックスの超伝導）：https://www.nobelprize.org/prizes/physics/年/summary/（curl で読んだ）
%%   ・高温超伝導のしくみの未解決：B. Keimer ら, "From quantum matter to high-temperature superconductivity in copper oxides", arXiv:1409.4673
%%       （要旨の "profound unresolved issues" を curl で読んだ。第4巻 D20 と同じ）
%%   ・ケイ素の禁制帯 1.12 eV（300 K）：Ioffe 研究所 NSM のデータ（www.ioffe.ru、作業環境から届かず検索の抜粋）と、BYU のクリーンルームの半導体定数の表
%%       （検索の抜粋）の2つで照らし合わせた。半導体の教科書（Sze）の値と同じ。
%%   ・石英ガラス（SiO₂）の禁制帯 約 9 eV（測り方で 8〜10 eV）：arXiv:1009.3873（非晶質 SiO₂ の反射率の論文の序論の要約）と、
%%       熱酸化膜の光伝導の測定（Suchow・Stemple、IBM の論文の記録）の検索の抜粋で照らし合わせた。本文は「約 9 eV」とし、はばを書いた。
%% 〔要確認〕（確かめられなかった値）
%%   ・銅のデバイ温度 343 K（固体物理の教科書 Kittel の表。教科書の表を引いた検索の抜粋1件。測り方で 315〜345 K ほどの違いがある）
%%   ・Nb-Ti の超伝導の臨界温度 約 9 K（製品の資料の抜粋で 9.2〜9.6 K。本文は「約 9.5 K とします」と書いた）
%%   ・GaAs の禁制帯 1.42 eV（じぶんでとこう1）：半導体の教科書の値（原-9 の組と合わせる）
\makeatletter
% 立方体（ななめから見た図）  \txicube{(左下手前)}{1辺}
\newcommand\txicube[2]{\begin{scope}[shift={#1}]
  \def\txiL{#2}
  \draw[bkInk!60, line width=0.5pt] (0,0) -- (\txiL,0) -- (\txiL,\txiL) -- (0,\txiL) -- cycle;
  \draw[bkInk!60, line width=0.5pt] (0,\txiL) -- ++(0.45*\txiL,0.3*\txiL) -- ++(\txiL,0) -- (\txiL,\txiL);
  \draw[bkInk!60, line width=0.5pt] (\txiL,0) -- ++(0.45*\txiL,0.3*\txiL) -- ++(0,\txiL);
  \draw[bkInk!40, line width=0.4pt, densely dashed] (0,0) -- ++(0.45*\txiL,0.3*\txiL) -- ++(\txiL,0) (0.45*\txiL,0.3*\txiL) -- ++(0,\txiL);
  \end{scope}}
% 立方体の中の点（u,v,w は 0〜1）  \txiatom{(左下手前)}{1辺}{u}{v}{w}{色}
\newcommand\txiatom[6]{\begin{scope}[shift={#1}]
  \filldraw[fill=#6, draw=bkInk, line width=0.4pt] ({#2*(#3+0.45*#5)},{#2*(#4+0.3*#5)}) circle (1.5);\end{scope}}
\tikzset{txi lab/.style={font=\sffamily\footnotesize, text=bkInk}, txi arr/.style={line width=0.9pt, -{Stealth[length=2.2mm]}},
  txi dim/.style={bkGray, line width=0.5pt, {Stealth[length=1.8mm]}-{Stealth[length=1.8mm]}}}
\makeatother

\phchapter{8}{固体の中の電子：\\エネルギーバンド}{金属は電気を通すのに、ガラスはなぜ通さない？}{\bklevelrange{2}{5}}{固体の中の電子：エネルギーバンド}
\bkmichishirube{4-5,4-6,4-13,5-14,7-8,9-3,熱-15,電-4,原-5,原-7}{原-9,原-10}

\begin{goal}
\begin{itemize}[label={\color{bkTheme}●}, leftmargin=1.6\zw, itemsep=0pt, topsep=1pt]
\item \bklevel{2}\enspace 金属・半導体・絶縁体の違いを、電子のエネルギーの帯（バンド）と、帯のすき間で説明できる。
\item \bklevel{3}\enspace ブラッグの条件で X線の回折の角を求め、格子定数から結晶の密度や原子の数を計算できる。
\item \bklevel{4}\enspace 3次元の箱で電子の席を数えて、金属のフェルミエネルギーとフェルミ速度を求められる。
\item \bklevel{4}\enspace フェルミ分布で、電子が席をうめる割合を求め、金属の電子の比熱が小さいわけを説明できる。
\item \bklevel{5}\enspace くし形の結晶の帯をプログラムで求め、平均自由行程・ウィーデマン・フランツの法則・デバイの $T^3$ 則を使える。
\end{itemize}
\end{goal}

\phkoushiki{%
\phk{金属のように電気を通す物、ガラスのように通さない物（絶縁体）、その間の半導体がある}{高校物理基礎}{ph:原-8:sec8}
\phk{X線の回折（ブラッグの条件）$2d\sin\theta=n\lambda$}{高校物理（発展）}{ph:原-8:bragg}
}

\begin{hatena}
\noindent\begin{minipage}[c]{0.56\linewidth}\raggedright
乾電池と豆電球の回路のあいだに、銅の線をはさむと、豆電球はつきます。ガラスの棒をはさむと、つきません。小学校でも調べたことです。

電気の通しやすさ（抵抗率の逆数）は、銅とガラスとで、何けたも違います（電-4）。銅の原子も、ガラスをつくるケイ素や酸素の原子も、同じように電子をもっているのに、なぜ、銅の電子は動けて、ガラスの電子は動けないのでしょう。
\end{minipage}\hfill
\begin{minipage}[c]{0.41\linewidth}\centering
\begin{tikzpicture}[x=0.86mm, y=0.86mm]
  \foreach \ox/\nm/\mat in {0/銅の線/1, 30/ガラスの棒/0} {
    \begin{scope}[shift={(\ox,0)}]
      \draw[bkInk, line width=0.7pt] (2,22) -- (2,30) -- (22,30) -- (22,22);
      \draw[bkInk, line width=0.7pt] (2,14) -- (2,4) -- (8,4) (16,4) -- (22,4) -- (22,14);
      \fill[bkGray!40] (0.5,14) rectangle (3.5,22); \draw[bkInk, line width=0.5pt] (0.5,14) rectangle (3.5,22);
      \filldraw[fill=bkYellow!40, draw=bkInk, line width=0.5pt] (22,18) circle (3.4);
      \ifnum\mat=1 \fill[bkOrange!80!black] (8,3.2) rectangle (16,4.8); \else \fill[bkBlue!25] (8,2.6) rectangle (16,5.4); \draw[bkInk!60, line width=0.4pt] (8,2.6) rectangle (16,5.4);\fi
      \node[font=\sffamily\footnotesize, anchor=north] at (12,1.6) {\nm};
    \end{scope}}
  \foreach \a in {0,45,...,315} {\draw[bkOrange!80!black, line width=0.6pt] ($(22,18)+(\a:4.4)$) -- ($(22,18)+(\a:6.2)$);}
  \node[font=\sffamily\footnotesize, anchor=south] at (12,31) {つく};
  \node[font=\sffamily\footnotesize, anchor=south] at (42,31) {つかない};
\end{tikzpicture}
\end{minipage}
\end{hatena}

\noindent 第2部（原-5〜原-7）では、1個の原子の中の電子の席を調べてきました。この章から始まる第3部では、原子が $10^{23}$ 個も集まった\textbf{固体}の中の電子を調べます。原子がきちんと並んだ固体の中では、電子の席は、とびとびの段ではなく、ぎっしりつまった帯（\textbf{エネルギーバンド}）になります。\bkidx{えねるぎーばんど}{エネルギーバンド}帯がどこまで電子でうまっているか、帯と帯のあいだのすき間がどれだけあるかで、金属・半導体・絶縁体が決まります。次の原-9・原-10 で学ぶダイオード・LED・トランジスタは、すべてこの章の「帯」の上に作られています。

\begin{bkfigure}
\begin{tikzpicture}[x=1mm, y=1mm]
  % 抵抗率の目もり（10 のべき -8〜16。1けた 4.6 mm）
  \fill[bkOrange!20] (0,0) rectangle (13.8,6);
  \fill[bkGreen!16] (13.8,0) rectangle (78.2,6);
  \fill[bkBlue!16] (78.2,0) rectangle (110.4,6);
  \draw[bkInk, line width=0.6pt] (0,0) rectangle (110.4,6);
  \node[txi lab] at (6.9,3) {金属};
  \node[txi lab] at (46,3) {半導体など（不純物や温度で大きく変わる）};
  \node[txi lab] at (94.3,3) {絶縁体};
  \draw[bkInk, line width=0.5pt] (0.0,0) -- (0.0,-1.5); \node[txi lab, anchor=north] at (0.0,-2) {$10^{-8}$};
  \draw[bkInk, line width=0.5pt] (18.4,0) -- (18.4,-1.5); \node[txi lab, anchor=north] at (18.4,-2) {$10^{-4}$};
  \draw[bkInk, line width=0.5pt] (36.8,0) -- (36.8,-1.5); \node[txi lab, anchor=north] at (36.8,-2) {$10^{0}$};
  \draw[bkInk, line width=0.5pt] (55.2,0) -- (55.2,-1.5); \node[txi lab, anchor=north] at (55.2,-2) {$10^{4}$};
  \draw[bkInk, line width=0.5pt] (73.6,0) -- (73.6,-1.5); \node[txi lab, anchor=north] at (73.6,-2) {$10^{8}$};
  \draw[bkInk, line width=0.5pt] (92.0,0) -- (92.0,-1.5); \node[txi lab, anchor=north] at (92.0,-2) {$10^{12}$};
  \draw[bkInk, line width=0.5pt] (110.4,0) -- (110.4,-1.5); \node[txi lab, anchor=north] at (110.4,-2) {$10^{16}$};
  \node[txi lab, anchor=west] at (113,-3.5) {\si{\ohm.m}};
  \draw[bkGray, line width=0.5pt] (1.06,6) -- (1.06,9); \fill[bkRed!75!black] (1.06,6) circle (0.9); \node[txi lab, anchor=south] at (1.06,9) {銅};
  \draw[bkGray, line width=0.5pt] (9.38,6) -- (9.38,16); \fill[bkRed!75!black] (9.38,6) circle (0.9); \node[txi lab, anchor=south] at (9.38,16) {ニクロム};
  \draw[bkGray, line width=0.5pt] (13.80,6) -- (13.80,9); \fill[bkRed!75!black] (13.80,6) circle (0.9); \node[txi lab, anchor=south] at (13.80,9) {黒鉛};
  \draw[bkGray, line width=0.5pt] (92.00,6) -- (92.00,9); \fill[bkRed!75!black] (92.00,6) circle (0.9); \node[txi lab, anchor=south] at (92.00,9) {ガラス};
  \draw[bkGray, line width=0.5pt] (110.40,6) -- (110.40,9); \fill[bkRed!75!black] (110.40,6) circle (0.9); \node[txi lab, anchor=south] at (110.40,9) {ポリエチレン};
\end{tikzpicture}
\bkfigcaption{いろいろな物の抵抗率の目安（電-4 の表の値。横は10のべきの目もりで、1目もりで $10^4$ 倍）。金属とプラスチックでは $10^{24}$ 倍も違う。ガラスは $10^{10}$〜$10^{14}$\,\si{\ohm.m} のはばのまん中に点を打った。}\label{ph:原-8:fig-rho}
\end{bkfigure}

\begin{yattemiyou}[結晶を見る・電気が通るか調べる]
\textbf{用意するもの}\par
食塩、ミョウバン（スーパーや薬局で買える）、虫めがね、黒い紙、乾電池（1.5 V）1個と電池ボックス、豆電球とソケット、みのむしクリップつきの導線3本、調べる物（鉛筆の芯、アルミはく、10円玉、ガラスのコップ、輪ゴム、紙）\par\vspace{2pt}
\begin{enumerate}[label={\bkcircled[bkGreen]{\arabic*}}, leftmargin=2.4\zw, itemsep=1pt, topsep=0pt]
\item 黒い紙の上に食塩とミョウバンのつぶを少し出し、虫めがねで形を見る。角の形と、平らな面を見つける。
\item 乾電池・豆電球・導線で回路を作り、とちゅうの2本の導線のあいだに、調べる物をはさむ。豆電球がつくかを見る。
\end{enumerate}
\par\vspace{2pt}\textbf{気をつけること}：乾電池は1個だけで使い、ショート（導線どうしを直接つなぐこと）させない。家のコンセントは、ぜったいに使わない。食塩やミョウバンは、口に入れない。
\end{yattemiyou}

\noindent やってみようの答え：食塩のつぶは、さいころのような立方体、ミョウバンは、角が8つの、ピラミッドを2つ合わせたような形です。どのつぶも、決まった角度の平らな面をもっています。中の原子（イオン）が、規則正しく並んでいるからです（①）。豆電球は、鉛筆の芯・アルミはく・10円玉ではつき、ガラス・輪ゴム・紙ではつきません。金属と黒鉛（鉛筆の芯）は電気を通し、ガラスやゴムは通しません。

\begin{omoidasou}
\bkrecall{原-5}{幅 $L$ の1次元の箱の電子のエネルギーは $E_n=\dfrac{n^2h^2}{8m_{\mathrm e}L^2}\ \ (n=1,2,\ldots)$。$L=\qty{1}{nm}$ で $E_1=\qty{0.376}{eV}$。波長は $\lambda=2L/n$。}
\bkrecall{原-7}{1つの席（量子数の組）には電子は1個まで（パウリの排他律）。原子が近づくと、電子の波は「結合」と「反結合」の2つの席に分かれる。}
\bkrecall{熱-15}{フェルミ分布 $\dfrac{1}{e^{(\varepsilon-\mu)/(k_{\mathrm B}T)}+1}$。1つの状態に入っている電子の平均の数。}
\bkrecall{電-4}{銅の自由電子の数は1 m$^3$ に $8.49\times10^{28}$ 個。ドルーデのモデルで、ぶつかるまでの平均の時間 $\tau=2.4\times10^{-14}$ s。導電率 $\sigma=ne^2\tau/m_{\mathrm e}$。}
\bkrecall{4-5}{半径 $R$ の球の体積は $\tfrac43\pi R^3$。}
\bkrecall{4-6}{直角三角形で、斜辺×$\sin\theta$ が、角 $\theta$ の向かいの辺。}
\end{omoidasou}

\noindent この章の記号：固体の本では、エネルギーを $E$ と書きます。熱-15 で $\varepsilon$ と書いた1個の粒子のエネルギーと同じものです。電子の数密度（1 m$^3$ あたりの数）を $n$ と書きます（第2部の主量子数とは別）。熱伝導率は、熱-11 では $k$ と書きましたが、この章では波数 $k$ とぶつかるので $\kappa$（カッパ）と書きます。

%% ======================== 物理の道具 ========================
\bksection{doguu}{物理の道具}

\phstage{me}{原子の並びを見る}{2}

\bksubsection{1}{結晶：原子の並び}
\label{ph:原-8:sec1}
多くの固体の中では、原子が、同じ形を繰り返して規則正しく並んでいます。このような固体を\textbf{結晶}といいます。\bkidx{けっしょう}{結晶}並び方は、小さな立方体などの「単位」を、上下・左右・前後にすき間なく積み重ねた形で表せます。図\ref{ph:原-8:fig-cubes}は、よく出てくる4つの並び方です。立方体の1辺の長さ $a$ を\textbf{格子定数}といいます。\bkidx{こうしていすう}{格子定数}

\begin{bkfigure}
\begin{tikzpicture}[x=1mm, y=1mm]
  % 単純立方
  \txicube{(0,0)}{18}
  \foreach \u/\v/\w in {0/0/0,1/0/0,0/1/0,1/1/0,0/0/1,1/0/1,0/1/1,1/1/1} {\txiatom{(0,0)}{18}{\u}{\v}{\w}{bkBlue!45}}
  \node[txi lab, anchor=north] at (13,-2) {単純立方：1個};
  % 体心立方
  \txicube{(33,0)}{18}
  \foreach \u/\v/\w in {0/0/0,1/0/0,0/1/0,1/1/0,0/0/1,1/0/1,0/1/1,1/1/1} {\txiatom{(33,0)}{18}{\u}{\v}{\w}{bkBlue!45}}
  \txiatom{(33,0)}{18}{0.5}{0.5}{0.5}{bkOrange!70}
  \node[txi lab, anchor=north, align=center] at (46,-2) {体心立方：2個\\（鉄）};
  % 面心立方
  \txicube{(66,0)}{18}
  \foreach \u/\v/\w in {0/0/0,1/0/0,0/1/0,1/1/0,0/0/1,1/0/1,0/1/1,1/1/1} {\txiatom{(66,0)}{18}{\u}{\v}{\w}{bkBlue!45}}
  \foreach \u/\v/\w in {0.5/0.5/0,0.5/0.5/1,0.5/0/0.5,0.5/1/0.5,0/0.5/0.5,1/0.5/0.5} {\txiatom{(66,0)}{18}{\u}{\v}{\w}{bkGreen!55}}
  \node[txi lab, anchor=north, align=center] at (79,-2) {面心立方：4個\\（銅・アルミニウム）};
  % ダイヤモンド形
  \txicube{(99,0)}{18}
  \foreach \u/\v/\w in {0/0/0,1/0/0,0/1/0,1/1/0,0/0/1,1/0/1,0/1/1,1/1/1} {\txiatom{(99,0)}{18}{\u}{\v}{\w}{bkBlue!45}}
  \foreach \u/\v/\w in {0.5/0.5/0,0.5/0.5/1,0.5/0/0.5,0.5/1/0.5,0/0.5/0.5,1/0.5/0.5} {\txiatom{(99,0)}{18}{\u}{\v}{\w}{bkGreen!55}}
  \foreach \u/\v/\w in {0.25/0.25/0.25,0.75/0.75/0.25,0.75/0.25/0.75,0.25/0.75/0.75} {\txiatom{(99,0)}{18}{\u}{\v}{\w}{bkRed!60}}
  \node[txi lab, anchor=north, align=center] at (112,-2) {ダイヤモンド形：8個\\（ケイ素・ダイヤモンド）};
\end{tikzpicture}
\bkfigcaption{結晶の4つの並び方（ななめから見た立方体。奥の辺は破線）。数は、立方体1つに入っている原子の数。かどの原子は8つの立方体に共有されるので $\tfrac18$ 個、面のまん中の原子は2つに共有されるので $\tfrac12$ 個、中の原子は1個と数える。}\label{ph:原-8:fig-cubes}
\end{bkfigure}

立方体1つの中の原子の数は、次のように数えます。
\begin{enumerate}[label={\bkcircled{\arabic*}}, leftmargin=2.4\zw, itemsep=1pt, topsep=2pt]
\item かどの原子は、周りの8つの立方体のかどを兼ねているので、1つの立方体には $\tfrac18$ 個ずつ。かどは8つで、$8\times\tfrac18=1$ 個。
\item 面のまん中の原子は、となりの立方体と半分ずつ。面は6つで、$6\times\tfrac12=3$ 個。
\item 立方体の中にある原子は、そのまま1個ずつ。
\end{enumerate}
単純立方は1個、体心立方は $1+1=2$ 個、面心立方は $1+3=4$ 個、ダイヤモンド形は、面心立方の4個に、中の4個を足して8個です。

銅は面心立方です。原-7 の例題8 で求めた銅の原子の数 $n=8.49\times10^{28}$ 個/m$^3$ を使うと、立方体1つに4個なので $a^3=4/n$、$a=\qty{0.361}{nm}$ です（じぶんでとこう3）。一番近い原子どうしの間かくは、面の対角線の半分で $a/\sqrt2=\qty{0.255}{nm}$ です。原子の大きさ（原-7 の電子の雲）と同じくらいで、原子どうしの雲が重なり合うほど近くに並んでいます。

\begin{phmijika}[食塩のつぶが、さいころの形なわけ]
食塩（塩化ナトリウム）の結晶は、ナトリウムのイオンと塩素のイオンが、立方体の形に交互に並んでいます。原子の並びの立方体が、何億個もそのまま積み重なって大きくなるので、目に見える大きさになっても、立方体の形を保ちます。雪の結晶が六角形なのも、水の分子の並び方が六角形の形をしているからです。
\end{phmijika}

\Needspace{12\baselineskip}\bksubsection{2}{X線で原子の並びを見る：ブラッグの条件}
\label{ph:原-8:sec2}
原子の間かくは $0.2$〜$0.3$ nm ほどで、光の顕微鏡では見えません（光の波長は 400〜700 nm：波-13）。そこで、波長が原子の間かくと同じくらいの X線を結晶に当てます。ラウエは、X線が結晶で回折して、決まった向きに強く散らばることを見つけました（1914年のノーベル物理学賞）。すぐにブラッグ父子は、この向きを、次のような簡単な考えで説明し、結晶の中の原子の並びを決める方法を作りました（1915年のノーベル物理学賞）。

\begin{bkfigure}
\begin{tikzpicture}[x=1mm, y=1mm]
  \foreach \yy in {-8,6,20} {\draw[bkGray, line width=0.4pt, densely dotted] (-4,\yy) -- (104,\yy);
    \foreach \xx in {0,10,...,100} {\fill[bkBlue!70!black] (\xx,\yy) circle (1.2);}}
  % 上の面で反射する X線（実線）と、下の面で反射する X線（破線）。角は 25°
  \draw[bkRed!80!black, txi arr] (13.75,36.9) -- (50,20);
  \draw[bkRed!80!black, txi arr] (50,20) -- (86.25,36.9);
  \draw[bkOrange!85!black, txi arr, dashed] (6.5,26.3) -- (50,6);
  \draw[bkOrange!85!black, txi arr, dashed] (50,6) -- (93.5,26.3);
  % 回り道
  \draw[bkInk, line width=0.5pt, densely dotted] (50,20) -- (44.64,8.50) (50,20) -- (55.36,8.50);
  \draw[bkGreen!50!black, line width=1.8pt] (44.64,8.50) -- (50,6) -- (55.36,8.50);
  \draw[bkGray, line width=0.5pt] (60,20) arc (0:25:10); \node[txi lab] at (64.0,22.6) {$\theta$};
  \draw[txi dim] (104,6) -- (104,20); \node[txi lab, anchor=west] at (105,13) {$d$};
  \node[txi lab, anchor=south] at (13.75,37.5) {X線};
  \node[txi lab, anchor=north] at (50,2.0) {太線の2本 $=2d\sin\theta$};
\end{tikzpicture}
\bkfigcaption{原子の面（点線）で反射する X線。上の面で反射する X線（実線）にくらべ、1つ下の面で反射する X線（破線）は、太線の部分（片側 $d\sin\theta$、合わせて $2d\sin\theta$）だけ長い道を通る。角 $\theta$ は、X線と面のあいだの角（図はかたむきを大きくかいた）。}\label{ph:原-8:fig-bragg}
\end{bkfigure}

結晶の中には、原子が並んだ平らな面が、間かく $d$ で何枚も重なっています（図\ref{ph:原-8:fig-bragg}）。X線は、それぞれの面で、鏡のように一部が反射します。上の面で反射した X線と、1つ下の面で反射した X線では、下の面のほうが $2d\sin\theta$ だけ長い道を通ります（4-6 の直角三角形）。この回り道が、波長のちょうど整数倍なら、すべての面で反射した X線の山と山がそろって、強め合います（波-13 の干渉）。
\pheq{2d\sin\theta=n\lambda\qquad(n=1,2,3,\ldots)}{m ＝ m}
\label{ph:原-8:bragg}これを\textbf{ブラッグの条件}といいます。\bkidx{ぶらっぐのじょうけん}{ブラッグの条件}

\begin{phyomi}[$2d\sin\theta=n\lambda$]
「面の間かくの2倍に $\sin\theta$ をかけた長さ（往きと帰りの回り道）が、波長の整数倍のときだけ、反射が強め合う」と読みます。ほかの角では、何千枚もの面から来る波が少しずつずれて、打ち消し合います。決まった波長の X線を当てて、強く反射する角 $\theta$ をはかれば、$d$ がわかります。$\sin\theta\le1$ なので、$\lambda\le2d$ でないと反射は起きません。可視光（数百 nm）は、原子の面の間かく（0.3 nm ほど）よりずっと長いので、結晶は光をブラッグ反射しません。
\end{phyomi}

\begin{phbaisuru}{面の間かく $d$}
$\sin\theta=n\lambda/(2d)$ は半分になります。角が小さいとき、$\sin\theta$ はほぼ $\theta$ に比例するので、強く反射する角もほぼ半分になります。間かくの大きい面ほど、小さい角で反射します。X線の回折の装置は、角を少しずつ変えながら反射の強さをはかり、強くなる角から、たくさんの面の間かくを求めて、原子の並びを決めます。
\end{phbaisuru}

\Needspace{10\baselineskip}\bksubsection{3}{ケイ素の結晶の密度と、原子の数}
\label{ph:原-8:sec3}
ケイ素の結晶は、ダイヤモンド形（立方体1つに8個）です。格子定数は、X線の回折でとても精密にはかられていて、$a=\qty{0.5431020511}{nm}$ です（CODATA 2022 の値）。これから、ケイ素の密度と、1 m$^3$ あたりの原子の数を計算できます（例題3）。
\pheq{\rho=\frac{8\,m_{\text{原子}}}{a^3},\qquad n=\frac{8}{a^3}}{kg/m$^3$ ＝ kg ÷ m$^3$}
\label{ph:原-8:density}ケイ素の原子1個の質量は $m_{\text{原子}}=\qty{28.09}{g/mol}\div N_{\mathrm A}$ です。計算すると、密度は \qty{2.329}{g/cm^3}、原子の数は $5.0\times10^{28}$ 個/m$^3$ になり、天びんではかった密度（\qty{2.330}{g/cm^3}）と、0.04\,\% の違いで合います。逆に、密度と格子定数をとても精密にはかれば、アボガドロ定数を決めることもできます。

\phstage{te}{電子の席を、固体の中で数える}{4}

\Needspace{12\baselineskip}\bksubsection{4}{自由電子を3次元の箱で数える：フェルミエネルギー}
\label{ph:原-8:sec4}
原-7 の例題8 で見たように、銅の原子は、いちばん外の $4s$ の電子を1個ずつ出して、金属の中を動き回らせています（自由電子）。この電子が、金属の中をどんなエネルギーで動いているかを考えます。

\begin{phmodel}[自由電子]
\phignore{原子の並び（イオン）が作る、でこぼこの位置エネルギー。電子どうしの反発。}
\phassume{金属のかたまりは1辺 $L$ の立方体の箱で、電子は箱の中を自由に動き、壁で閉じこめられる（原-5 の箱の3次元版）。}
\phlimit{電子の波長が、原子の間かくの2倍に近いとき（ブラッグ反射：⑦）。$d$ の電子が主役の金属（鉄・ニッケルなど）。}
\end{phmodel}

\phboxhead{3次元の箱の席}
原-5 ③の1次元の箱では、エネルギーは $E=\dfrac{h^2}{8m_{\mathrm e}L^2}\,n^2$ でした。3次元の立方体の箱では、$x$・$y$・$z$ の3つの向きに、それぞれ定在波の腹の数 $n_x$・$n_y$・$n_z$（どれも $1,2,3,\ldots$）があり、エネルギーはそれぞれの向きの分の和になります。
\pheq{E=\frac{h^2}{8m_{\mathrm e}L^2}\bigl(n_x^{\,2}+n_y^{\,2}+n_z^{\,2}\bigr)}{J ＝ (J·s)$^2$ ÷ (kg·m$^2$)}
3つの数の組 $(n_x,n_y,n_z)$ に、スピンの向き2つをかけたものが、電子の席です。

\phboxhead{席を、球の体積で数える}
$(n_x,n_y,n_z)$ を、3次元の座標の点と考えます。点は、1辺1の立方体のかどに、1つずつあります。エネルギーが $E$ 以下の席は、$n_x^{\,2}+n_y^{\,2}+n_z^{\,2}\le R^2$、つまり原点からの距離が $R$ 以下の点です（$R^2=8m_{\mathrm e}L^2E/h^2$）。$n_x,n_y,n_z$ はどれも正なので、点は、半径 $R$ の球の $\tfrac18$（どの座標も正の部分）に入っています。点1つが体積1を受け持つので、$R$ が大きければ、点の数は体積で数えられます（4-5）。
\[ \text{点の数}\approx\frac18\times\frac43\pi R^3=\frac{\pi}{6}R^3 \]
図\ref{ph:原-8:fig-count}の表のように、$R$ が大きいほど、この数え方は正確になります。本物の金属では $R$ は1億くらいなので、まったく正確です。

\begin{bkfigure}
\begin{minipage}[c]{0.46\linewidth}\centering
\begin{tikzpicture}[x=1mm, y=1mm]
  \draw[bkInk, txi arr] (0,0) -- (48,0); \draw[bkInk, txi arr] (0,0) -- (0,48);
  \node[txi lab, anchor=north] at (47,-1) {$n_x$}; \node[txi lab, anchor=east] at (-1,47) {$n_y$};
  \fill[bkBlue!10] (0,0) -- (40,0) arc (0:90:40) -- cycle;
  \draw[bkBlue!80!black, line width=0.9pt] (40,0) arc (0:90:40);
  \foreach \i in {1,...,9} {\foreach \j in {1,...,9} {
    \pgfmathparse{(\i*\i+\j*\j<=64)?1:0}
    \ifnum\pgfmathresult=1 \fill[bkOrange!85!black] (\i*5,\j*5) circle (1.0); \else \draw[bkGray, line width=0.4pt] (\i*5,\j*5) circle (1.0);\fi}}
\end{tikzpicture}
\end{minipage}\hfill
\begin{minipage}[c]{0.5\linewidth}\centering
\begin{bktable}{colspec={Q[c,m,wd=10mm] Q[c,m,wd=22mm] Q[c,m,wd=20mm]}, rowsep=1pt, row{1}={bg=bkTheme!10}, cells={font=\small}}
$R$ & 数えた点の数 & $\dfrac{\pi}{6}R^3$ \\
5 & 38 & 65 \\
10 & 410 & 524 \\
20 & \num{3719} & \num{4189} \\
40 & \num{31616} & \num{33510} \\
\end{bktable}
\end{minipage}
\bkfigcaption{左：2次元にして見た、席の点（2次元では円の $\tfrac14$。円の半径は $R=8$）。半径 $R$ の内側の点（ぬりつぶした点）がエネルギー $E$ 以下の席。右：3次元で、半径 $R$ の球の $\tfrac18$ の中の点を、Python で1つずつ数えた数と、体積の式の比較。$R$ が大きいほど近づく（足りない分は、球の表面のあたりの点）。}\label{ph:原-8:fig-count}
\end{bkfigure}

\phboxhead{いちばん上の席：フェルミエネルギー}
金属の中に電子が $N$ 個あります。パウリの排他律（原-7 ⑦）のために、電子は1つの席に1個ずつ、エネルギーの低い席から順に入っていきます。いちばん温度の低い（0 K の）とき、$N$ 個の電子は、半径 $R_{\mathrm F}$ の $\tfrac18$ 球の中の席を、ちょうどうめます。スピンで席は2倍なので、
\[ N=2\times\frac{\pi}{6}R_{\mathrm F}^{\,3}\qquad\text{より}\qquad R_{\mathrm F}=\Bigl(\frac{3N}{\pi}\Bigr)^{1/3} \]
いちばん上の電子のエネルギーは、$E=\dfrac{h^2}{8m_{\mathrm e}L^2}R_{\mathrm F}^{\,2}$ に入れて、
\pheq{E_{\mathrm F}=\frac{h^2}{8m_{\mathrm e}}\Bigl(\frac{3n}{\pi}\Bigr)^{2/3}=\frac{\hbar^2}{2m_{\mathrm e}}\bigl(3\pi^2n\bigr)^{2/3}}{J ＝ (J·s)$^2$ ÷ kg × m$^{-2}$}
\label{ph:原-8:EF}です（$n=N/L^3$ は電子の数密度。2つの形は、$h=2\pi\hbar$ を入れると同じになります）。この、0 K でいちばん上までうまった席のエネルギーを\textbf{フェルミエネルギー}といいます。\bkidx{ふぇるみえねるぎー}{フェルミエネルギー}箱の大きさ $L$ が消えて、数密度 $n$ だけで決まるのが大事な点です。

\begin{phyomi}[$E_{\mathrm F}=\dfrac{\hbar^2}{2m_{\mathrm e}}(3\pi^2n)^{2/3}$]
「電子がぎっしりつまっているほど、いちばん上の電子のエネルギーは高い」と読みます。電子は同じ席に入れないので、たくさんの電子を入れると、上の席まで使うしかないのです。銅（$n=8.49\times10^{28}$ /m$^3$）では $E_{\mathrm F}=\qty{7.0}{eV}$ です（例題4）。原-5 の1 nm の箱の電子（0.376 eV）より、ずっと大きなエネルギーです。
\end{phyomi}

いちばん上の電子の速さ（\textbf{フェルミ速度}）は、$\tfrac12m_{\mathrm e}v_{\mathrm F}^{\,2}=E_{\mathrm F}$ から求めます。銅では $v_{\mathrm F}=\qty{1.57e6}{m/s}$ で、1秒に約 \qty{1570}{km}です。電-4 で「自由電子は、ばらばらな向きに1秒に千 km 以上の速さで動き回っている」と書いたのは、この速さです。温度で決まる速さではなく、\textbf{0 K でも}、電子は排他律のために、この速さで動き回っているのです。$E_{\mathrm F}$ を温度に直した $T_{\mathrm F}=E_{\mathrm F}/k_{\mathrm B}$（\textbf{フェルミ温度}）は、銅で約 \qty{8.2e4}{K}もあります。

\begin{phbaisuru}[8倍]{電子の数密度 $n$}
$E_{\mathrm F}\propto n^{2/3}$ なので、$8^{2/3}=4$ 倍になります。電子が8倍つまると、1辺の席の数（$R_{\mathrm F}$）が2倍になり、エネルギーはその2乗の4倍です。アルミニウムは原子1個から電子を3個出すので、$n$ は銅の約2.1倍で、$E_{\mathrm F}$ は約 \qty{11.7}{eV}です。ナトリウムは、原子の間かくが広くて $n$ が小さいので、約 \qty{3.2}{eV}です。
\end{phbaisuru}

\Needspace{12\baselineskip}\bksubsection{5}{フェルミ分布：温度があると、ふちがぼやける}
\label{ph:原-8:sec5}
0 K では、$E_{\mathrm F}$ より下の席はすべてうまり、上の席はすべて空いています。温度が上がると、熱の動きで、$E_{\mathrm F}$ のすぐ下の電子の一部が、すぐ上の空いた席に上がります。エネルギー $E$ の席に電子がいる確率（平均の数）は、熱-15 ⑨の\textbf{フェルミ分布}
\pheq{f(E)=\frac{1}{e^{(E-E_{\mathrm F})/(k_{\mathrm B}T)}+1}}{1 ＝ 1}
\label{ph:原-8:fermi}です。\bkidx{ふぇるみぶんぷ}{フェルミ分布}熱-15 では、1個の粒子のエネルギーを $\varepsilon$、化学ポテンシャルを $\mu$ と書きました。固体の本では、これを $E$ と $E_{\mathrm F}$ と書きます（$\mu$ は温度でわずかに変わりますが、金属ではほぼ $E_{\mathrm F}$ と同じです）。

\begin{bkfigure}
\begin{tikzpicture}
\begin{axis}[bkplot, axis lines=left, width=110mm, height=48mm, xmin=-0.4, xmax=0.4, ymin=0, ymax=1.15, xtick={-0.4,-0.2,0,0.2,0.4}, xticklabels={$-0.4$,$-0.2$,0,0.2,0.4}, ytick={0,0.5,1}, xlabel={$E-E_{\mathrm F}$（eV）}, ylabel={$f(E)$}]
  \addplot[bkcurve, domain=-0.4:0.4, samples=200] {1/(exp(x/0.025852)+1)};
  \addplot[bkcurve2, dashed, domain=-0.4:0.4, samples=200] {1/(exp(x/0.0861733)+1)};
  \draw[bkInk, line width=0.8pt, densely dotted] (axis cs:-0.4,1) -- (axis cs:0,1) -- (axis cs:0,0) -- (axis cs:0.4,0);
  \node[font=\sffamily\footnotesize, anchor=west] at (axis cs:0.04,0.9) {0 K（点線）：段};
  \node[font=\sffamily\footnotesize, anchor=west] at (axis cs:0.06,0.42) {300 K（実線）};
  \node[font=\sffamily\footnotesize, anchor=west] at (axis cs:0.17,0.2) {1000 K（破線）};
\end{axis}
\end{tikzpicture}
\bkfigcaption{フェルミ分布。0 K では $E_{\mathrm F}$ で0と1が切りかわる段。温度があると、$E_{\mathrm F}$ の周りの $k_{\mathrm B}T$ の数倍のはば（300 K で 0.1 eV ほど）だけ、段がなだらかになる。$E=E_{\mathrm F}$ では、どの温度でも $f=\tfrac12$。}\label{ph:原-8:fig-fermi}
\end{bkfigure}

\begin{phyomi}[$f(E)=1/(e^{(E-E_{\mathrm F})/(k_{\mathrm B}T)}+1)$]
「$E_{\mathrm F}$ よりずっと下の席はほぼ必ずうまり（$f\approx1$）、ずっと上の席はほぼ空いている（$f\approx0$）。変わり目のはばは $k_{\mathrm B}T$ の数倍」と読みます。$E-E_{\mathrm F}$ が $k_{\mathrm B}T$ の数倍より大きいと、分母の1が無視できて、$f\approx e^{-(E-E_{\mathrm F})/(k_{\mathrm B}T)}$ のボルツマン因子（熱-15）になります。原-9 では、半導体の電子の数を、この形で数えます。
\end{phyomi}

\begin{phbaisuru}{温度 $T$}
図\ref{ph:原-8:fig-fermi}のように、段のなだらかな部分のはばが2倍になります。たとえば、$f$ が 0.9 から 0.1 に下がるはばは $2k_{\mathrm B}T\ln9\approx4.4k_{\mathrm B}T$ で、300 K で 0.11 eV、600 K で 0.23 eV です。それでも、銅の $E_{\mathrm F}=\qty{7.0}{eV}$ にくらべると、ほんのわずかです。
\end{phbaisuru}

\Needspace{12\baselineskip}\bksubsection{6}{金属の電子の比熱が、小さいわけ}
\label{ph:原-8:sec6}
電子が気体の分子のように自由に動くなら、熱-14 の等分配の法則で、電子1個は $\tfrac32k_{\mathrm B}T$ のエネルギーをもち、1 mol の電子の比熱は $\tfrac32R\approx\qty{12.5}{J/(mol.K)}$ のはずです。銅は原子1個に電子1個なので、原子のゆれの比熱 $3R$（熱-15 のデュロン・プティの値）の半分も、上乗せされるはずです。ところが、はかった金属の比熱は、ほぼ $3R$ だけで、電子の分はほとんど見えません。19世紀の終わりから、ドルーデのモデル（電-4）を悩ませた大きななぞでした。

フェルミ分布が答えを出します。温度 $T$ で熱のエネルギーを受け取れるのは、$E_{\mathrm F}$ の周りの $k_{\mathrm B}T$ くらいのはばにいる電子だけです。それより深い電子は、$k_{\mathrm B}T$ くらいのエネルギーをもらっても、上の席がすべてうまっていて、移れる空いた席がないのです（パウリの排他律）。熱を受け取れる電子の割合は、およそ $k_{\mathrm B}T/E_{\mathrm F}=T/T_{\mathrm F}$ で、300 K の銅では $300/\num{82000}\approx0.004$ にすぎません。くわしく計算すると（ゾンマーフェルト）、電子の比熱は
\pheq{C_{\text{電子}}=\frac{\pi^2}{2}\,Nk_{\mathrm B}\,\frac{T}{T_{\mathrm F}}}{J/K ＝ J/K}
\label{ph:原-8:Cel}で、等分配の $\tfrac32Nk_{\mathrm B}$ の $\dfrac{\pi^2}{3}\cdot\dfrac{T}{T_{\mathrm F}}$ 倍、300 K の銅で約 0.012 倍です（例題6）。

\begin{phmijika}[満員電車の、奥の人は動けない]
ぎっしり満員の電車を思いうかべましょう。ドアの近くの人は、駅に着くと、空いたホームに出たり、また入ったりできます。でも、まん中の奥にいる人は、周りがすべてうまっていて、動く場所がありません。金属の中の電子も同じで、熱の出入りに関われるのは、$E_{\mathrm F}$ の近くの「ドアの近く」の電子だけなのです。
\end{phmijika}

\Needspace{14\baselineskip}\bksubsection{7}{結晶の中の波は、ブラッグ反射して止まる：帯のすき間}
\label{ph:原-8:sec7}
④の自由電子のモデルでは、電子は原子の並びを感じませんでした。本物の結晶の中では、電子は、規則正しく並んだイオン（原子核と内側の電子）の作る、でこぼこの位置エネルギーの中を動きます。そのでこぼこが、電子の席に何をもたらすかを、2つの見方で考えます。

\phboxhead{見方1：原子の席が集まって、帯になる}
原-7 ⑪で、2つの原子を近づけると、1つの席が「結合」と「反結合」の2つに分かれました。原子を4個、8個と並べると、席は4つ、8つに分かれ、たがいに少しずつ違うエネルギーになります。$N$ 個（$10^{23}$ 個）並べると、$N$ 個の席が、決まったはばのエネルギーの範囲にぎっしりつまります。これが\textbf{帯}（バンド）です。帯のはばは、となりの原子の電子の雲がどれだけ重なるかで決まり、数 eV ほどです。1つの帯には、原子1個の1つの席から生まれた $N$ 個の席があり、スピンの2つをかけて、$2N$ 個の電子が入れます。違う席（$3s$ と $3p$ など）から生まれた帯のあいだには、どの席もない\textbf{すき間}（禁制帯）ができることがあります。\bkidx{きんせいたい}{禁制帯}

\begin{bkfigure}
\begin{tikzpicture}[x=1mm, y=1mm]
  \draw[bkBlue!80!black, line width=0.7pt] (0,8) -- (12,8);
  \draw[bkOrange!85!black, line width=0.7pt] (0,26) -- (12,26);
  \node[txi lab, anchor=north] at (6,-0.5) {1個};
  \foreach \xx/\nn/\lab in {22/2/2個, 44/4/4個, 66/8/8個} {
    \pgfmathsetmacro\wA{3.2*(1-1/\nn)}
    \pgfmathsetmacro\wB{5.2*(1-1/\nn)}
    \foreach \k in {1,...,\nn} {
      \pgfmathsetmacro\yA{8-\wA+2*\wA*(\k-1)/(\nn-1)}
      \pgfmathsetmacro\yB{26-\wB+2*\wB*(\k-1)/(\nn-1)}
      \draw[bkBlue!80!black, line width=0.7pt] (\xx,\yA) -- (\xx+12,\yA);
      \draw[bkOrange!85!black, line width=0.7pt] (\xx,\yB) -- (\xx+12,\yB);}
    \node[txi lab, anchor=north] at (\xx+6,-0.5) {\lab};}
  % たくさん
  \fill[bkBlue!25] (88,4.8) rectangle (100,11.2); \draw[bkBlue!80!black, line width=0.7pt] (88,4.8) rectangle (100,11.2);
  \fill[bkOrange!25] (88,20.8) rectangle (100,31.2); \draw[bkOrange!85!black, line width=0.7pt] (88,20.8) rectangle (100,31.2);
  \node[txi lab, anchor=north] at (94,-0.5) {$N$ 個};
  \draw[txi dim] (104,11.2) -- (104,20.8); \node[txi lab, anchor=west] at (105,16) {すき間};
  \node[txi lab, anchor=west] at (102,26) {帯};
  \node[txi lab, anchor=west] at (102,8) {帯};
  \draw[bkInk, txi arr] (-5,0) -- (-5,34); \node[txi lab, anchor=south] at (-5,34.5) {エネルギー};
\end{tikzpicture}
\bkfigcaption{原子をたくさん並べると、原子1個の1つの席（左はし）が、原子の数だけの席に分かれ、ぎっしりつまった帯になる。原子1個の下の席と上の席から、それぞれ帯ができ、あいだにすき間が残る。}\label{ph:原-8:fig-split}
\end{bkfigure}

\phboxhead{見方2：自由な電子の波が、ブラッグ反射する}
反対に、④の自由電子から出発します。間かく $a$ でイオンが並んだ1次元の結晶の中を、電子の波が進むとします。それぞれのイオンは、電子の波をほんの少しずつはね返します。ふつうは、はね返った波どうしの山と谷がずれて、打ち消し合います。ところが、電子の波長が
\[ \lambda=2a\qquad\Bigl(\text{波数}\ k=\frac{2\pi}{\lambda}=\frac{\pi}{a}\Bigr) \]
のときは、となりのイオンではね返る波は、往きと帰りで $2a$、ちょうど1波長だけ長い道を通るので、すべての反射がそろって強め合います。②の X線のブラッグの条件（$2d\sin\theta=n\lambda$ で $\theta=90^\circ$、$n=1$）と同じです。電子の波は完全にはね返され、進めなくなります。進む波とはね返った波が重なって、定在波になるのです（原-5 ①の弦と同じ）。

\begin{bkfigure}
\begin{tikzpicture}
\begin{axis}[bkplot, axis lines=left, name=pk, width=54mm, height=54mm, xmin=-1.6, xmax=1.6, ymin=0, ymax=2.6, xtick={-1,0,1}, xticklabels={$-\frac{\pi}{a}$,0,$\frac{\pi}{a}$}, ytick=\empty, xlabel={波数 $k$}, ylabel={$E$}]
  \addplot[bkGray, densely dotted, line width=0.8pt, domain=-1.55:1.55] {x^2};
  \addplot[bkcurve, domain=-1:1, samples=80] {x^2 - 0.18*(x^2)^6};
  \addplot[bkcurve, domain=1:1.5, samples=40] {x^2 + 0.18 + 0.18*exp(-12*(x-1))};
  \addplot[bkcurve, domain=-1.5:-1, samples=40] {x^2 + 0.18 + 0.18*exp(12*(x+1))};
  \draw[bkInk, line width=0.6pt, {Stealth[length=1.5mm]}-{Stealth[length=1.5mm]}] (axis cs:1.08,0.82) -- (axis cs:1.08,1.36);
  \node[font=\sffamily\footnotesize, anchor=west] at (axis cs:1.1,1.09) {すき間};
\end{axis}
\begin{scope}[x=1mm, y=1mm, shift={(66,4)}]
  \foreach \xx in {0,12,24,36,48} {\fill[bkRed!70!black] (\xx,0) circle (1.3); \fill[bkRed!70!black] (\xx,24) circle (1.3);}
  \draw[bkBlue!80!black, line width=1.0pt, domain=-2:50, samples=80, smooth] plot (\x,{1.6+9*cos(\x*15)^2});
  \draw[bkOrange!85!black, line width=1.0pt, dashed, domain=-2:50, samples=80, smooth] plot (\x,{25.6+9*sin(\x*15)^2});
  \node[txi lab, anchor=south] at (24,11.5) {$\cos^2$（実線）：イオンの上 → 低い};
  \node[txi lab, anchor=south] at (24,35.5) {$\sin^2$（破線）：イオンのあいだ → 高い};
  \draw[txi dim] (0,-3) -- (12,-3); \node[txi lab, anchor=north] at (6,-3.5) {$a$};
\end{scope}
\end{tikzpicture}
\bkfigcaption{左：自由電子のエネルギー $E\propto k^2$（点線）は、$k=\pm\pi/a$ でブラッグ反射のために切れて、すき間ができる（実線）。右：$k=\pi/a$ の2つの定在波の、電子が見つかる確率 $|\psi|^2$。点はイオンの位置。}\label{ph:原-8:fig-gap}
\end{bkfigure}

定在波は2通りあります。1つは、山（$|\psi|^2$ の大きい所）がイオンの上にくる $|\psi|^2\propto\cos^2(\pi x/a)$ で、電子が正のイオンの近くに多いので、位置エネルギーが低くなります。もう1つは、山がイオンとイオンのあいだにくる $|\psi|^2\propto\sin^2(\pi x/a)$ で、位置エネルギーが高くなります。同じ波長なのに、エネルギーが2つに分かれるのです。そのあいだのエネルギーをもつ、進む波はありません。これが帯のすき間です。すき間ができる所の自由電子のエネルギーは、
\pheq{E=\frac{\hbar^2}{2m_{\mathrm e}}\Bigl(\frac{\pi}{a}\Bigr)^2=\frac{h^2}{8m_{\mathrm e}a^2}}{J ＝ (J·s)$^2$ ÷ (kg·m$^2$)}
\label{ph:原-8:gapE}で、原-5 の幅 $a$ の箱のいちばん低いエネルギーと同じ形です。$a=\qty{0.25}{nm}$ なら約 \qty{6.0}{eV}です（例題7）。$k=2\pi/a$、$3\pi/a$、…（$2a=n\lambda$）でも、同じようにすき間ができます。2つの見方は、同じ帯とすき間を、原子の側と波の側から見たものです。

\Needspace{14\baselineskip}\bksubsection{8}{導体・半導体・絶縁体}
\label{ph:原-8:sec8}
電場をかけて電流が流れるには、電子が、電場の向きに少しだけ速くなる、つまり、少しだけエネルギーの高い、すぐ上の席に移れなければなりません。

\begin{enumerate}[label={\bkcircled{\arabic*}}, leftmargin=2.4\zw, itemsep=1pt, topsep=2pt]
\item \textbf{帯が電子でぎっしりうまっているとき}：どの席もうまっているので、電子は移る先がありません。右向きに動く電子と左向きに動く電子がちょうど同じ数だけいて、その組み合わせを変えられないので、電流は流れません。
\item \textbf{帯が途中までうまっているとき}：$E_{\mathrm F}$ のすぐ上に空いた席がいくらでもあるので、弱い電場でも電子は少し速い席に移れます。電流が流れます。
\end{enumerate}

\begin{bkfigure}
\begin{tikzpicture}[x=1mm, y=1mm]
  \foreach \xx/\nm/\lo/\hi/\fil in {0/金属（導体）/0/0/1, 42/半導体/14/20/0, 84/絶縁体/14/34/0} {
    \begin{scope}[shift={(\xx,0)}]
      \ifnum\fil=1
        \fill[bkBlue!30] (0,4) rectangle (26,16);
        \fill[pattern=north east lines, pattern color=bkBlue!80!black] (0,4) rectangle (26,16);
        \draw[bkInk, line width=0.6pt] (0,4) rectangle (26,26);
        \draw[bkRed!80!black, line width=1.0pt] (-1,16) -- (27,16); \node[txi lab, anchor=west] at (27,16) {$E_{\mathrm F}$};
      \else
        \fill[pattern=north east lines, pattern color=bkBlue!80!black] (0,4) rectangle (26,\lo);
        \draw[bkInk, line width=0.6pt] (0,4) rectangle (26,\lo);
        \draw[bkInk, line width=0.6pt] (0,\hi) rectangle (26,\hi+8);
        \draw[txi dim] (13,\lo) -- (13,\hi);
        \node[txi lab, anchor=west] at (14,{(\lo+\hi)/2}) {$E_{\mathrm g}$};
      \fi
      \node[txi lab, anchor=north] at (13,2.5) {\nm};
    \end{scope}}
\end{tikzpicture}
\bkfigcaption{帯のうまり方（斜線の部分が電子でうまった席、白い部分が空いた席）。金属は帯が途中までうまっていて、$E_{\mathrm F}$ のすぐ上に空いた席がある。半導体と絶縁体は、下の帯（価電子帯）がぎっしりうまり、すき間 $E_{\mathrm g}$ をはさんで、上の帯（伝導帯）が空いている。すき間が小さいのが半導体、大きいのが絶縁体。}\label{ph:原-8:fig-3bands}
\end{bkfigure}

ナトリウムや銅は、原子1個から $s$ の電子を1個出します。$s$ の席から生まれた帯は $2N$ 個の電子が入れるので、$N$ 個の電子では\textbf{半分}までしかうまりません。だから金属です。マグネシウムは原子1個から2個出すので、$s$ の帯がちょうどうまりそうですが、$3s$ と $3p$ の帯が重なっていて、空いた席がすぐ上にあるので、やはり金属です。

ケイ素やダイヤモンドは、原子1個が4個の電子を出してとなりの4個の原子と共有結合を作ります（原-7 ⑪）。結合の席から生まれた下の帯（\textbf{価電子帯}）は、ちょうどぎっしりうまり、反結合の席から生まれた上の帯（\textbf{伝導帯}）は空いています。\bkidx{かでんしたい}{価電子帯}\bkidx{でんどうたい}{伝導帯}そのあいだのすき間の大きさを\textbf{禁制帯幅} $E_{\mathrm g}$ といいます。ケイ素は $E_{\mathrm g}=\qty{1.12}{eV}$、ガラスの主な成分の二酸化ケイ素は約 \qty{9}{eV}（はかり方によって 8〜10 eV）です。すき間が小さく、室温で少しは電子が上の帯へ上がるものを\textbf{半導体}、すき間が大きく、ほとんど上がらないものを\textbf{絶縁体}といいます。\bkidx{はんどうたい}{半導体}\bkidx{ぜつえんたい}{絶縁体}

\phboxhead{上の帯へ上がれる電子の割合}
不純物のない半導体では、$E_{\mathrm F}$ はすき間のちょうどまん中あたりにあります（上の帯の電子の数と、下の帯の空いた席の数が等しくなるように決まるから：原-9 ③）。上の帯の底は、$E_{\mathrm F}$ より $E_{\mathrm g}/2$ だけ上です。そこの席に電子がいる確率は、⑤のフェルミ分布のすそ（ボルツマン因子）で、
\pheq{f\approx e^{-E_{\mathrm g}/(2k_{\mathrm B}T)}}{1 ＝ 1}
\label{ph:原-8:egfactor}です。$E_{\mathrm g}$ が $k_{\mathrm B}T$ の何十倍もあるので、この数はとても小さく、しかも $E_{\mathrm g}$ によって何十けたも変わります（例題1）。半導体の電子の数をきちんと数えるのは、原-9 ③です。

\begin{phbaisuru}{禁制帯幅 $E_{\mathrm g}$}
$e^{-2E_{\mathrm g}/(2k_{\mathrm B}T)}=\bigl(e^{-E_{\mathrm g}/(2k_{\mathrm B}T)}\bigr)^2$ なので、割合は\textbf{2乗}になります。ケイ素（1.12 eV）の $4.1\times10^{-10}$ が、2.24 eV では $1.7\times10^{-19}$ です（300 K）。「2倍」ではなく「2乗」で、何けたも小さくなるのが、指数の式の特徴です。反対に、温度を上げると割合は急に増え、ケイ素では 300 K から 350 K で約22倍になります。半導体の抵抗が、温度を上げると下がるのはこのためです（電-4 のサーミスター）。
\end{phbaisuru}

\Needspace{12\baselineskip}\bksubsection{9}{はてな？の答え}
\label{ph:原-8:sec9}
\phboxhead{\bkhatenaref の答え}
\label{ph:原-8:hatena}銅の原子は $4s$ の電子を1個ずつ出し、その帯は半分までしかうまっていません。$E_{\mathrm F}$（7.0 eV）のすぐ上に空いた席がぎっしりあるので、電池の弱い電場でも、たくさんの電子が少し速い席に移り、電流が流れます。

ガラスの中の電子は、ケイ素と酸素の結合の席から生まれた価電子帯をぎっしりうめていて、すぐ上には席がありません。上の帯までは約 9 eV のすき間があり、室温の熱のエネルギー（$k_{\mathrm B}T=\qty{0.026}{eV}$）では、上がれる電子の割合は $e^{-9/(2\times0.0259)}\approx3\times10^{-76}$ です。1 m$^3$ のガラスの中の電子は、全部合わせても $10^{30}$ 個ほどなので、上の帯の電子は1個もいないのと同じです。だから、ガラスは電気を通しません。（ガラスがほんの少しだけ電気を通すのは、ガラスの中のナトリウムなどのイオンが、ゆっくり動くためです。）

ケイ素の割合は $4.1\times10^{-10}$ で、ガラスより約 $10^{66}$ 倍も大きく、室温でもいくらか電気を通します（半導体）。電-4 で、金属とプラスチックの抵抗率が $10^{24}$ 倍も違うのを見ました。その違いのもとは、帯が途中までうまっているか、すき間をはさんでぎっしりうまっているか、そして、そのすき間の大きさです。

\phkokomade

\Needspace{9\baselineskip}\phstage{shiki}{帯を式とプログラムで調べる}{5}

\bksubsection{10}{平均自由行程と、ウィーデマン・フランツの法則}
\label{ph:原-8:sec10}
\phboxhead{電子は、イオンにぶつかっていない}
電-4 のドルーデのモデルでは、銅の電子は、平均 $\tau=\qty{2.4e-14}{s}$ ごとに何かにぶつかっていました。ぶつかるまでに進む距離（\textbf{平均自由行程} $\ell$）は、電子の速さに $\tau$ をかけたものです。⑥で見たように、電気を運ぶのは $E_{\mathrm F}$ の近くの電子で、その速さは $v_{\mathrm F}$ なので、
\pheq{\ell=v_{\mathrm F}\tau}{m ＝ m/s × s}
\label{ph:原-8:mfp}銅では $1.57\times10^6\times2.4\times10^{-14}\approx\qty{38}{nm}$、原子の間かく（0.256 nm）の約150倍です（例題9）。電子は、イオンを1個1個よけて、ずっと遠くまで進んでいるのです。

これは、電子が波だからです。完全に規則正しく並んだイオンの中では、電子の波は、ブラッグ反射の条件（⑦）に合う波を除いて、何にもぶつからずにまっすぐ進み続けます（ブロッホの定理：専門書の囲み）。電子をはね返すのは、並びの\textbf{乱れ}、つまり、熱でゆれる原子と、まじった不純物や結晶の欠陥です。温度が上がると、原子のゆれが大きくなって $\tau$ が短くなり、抵抗率が大きくなります（電-4 ④）。

\phboxhead{熱を運ぶのも、同じ電子}
金属は、電気をよく通すものほど、熱もよく伝えます（銀・銅・金・アルミニウムの順）。熱も、$E_{\mathrm F}$ の近くの電子が運んでいるからです。熱-14 ⑧の気体の熱伝導率の式（$\tfrac13\times$比熱$\times$速さ$\times$平均自由行程）を、金属の電子に使います。1 m$^3$ あたりの電子の比熱は⑥の式で $\tfrac{\pi^2}{2}nk_{\mathrm B}\dfrac{T}{T_{\mathrm F}}$、速さは $v_{\mathrm F}$、平均自由行程は $v_{\mathrm F}\tau$ です。$k_{\mathrm B}T_{\mathrm F}=E_{\mathrm F}=\tfrac12m_{\mathrm e}v_{\mathrm F}^{\,2}$ を使うと、
\[ \kappa=\frac13\cdot\frac{\pi^2}{2}nk_{\mathrm B}\frac{T}{T_{\mathrm F}}\cdot v_{\mathrm F}\cdot v_{\mathrm F}\tau=\frac{\pi^2nk_{\mathrm B}^{\,2}T\tau}{3m_{\mathrm e}} \]
電-4 の導電率 $\sigma=\dfrac{ne^2\tau}{m_{\mathrm e}}$ で割ると、$n$ も $\tau$ も $m_{\mathrm e}$ も消えて、
\pheq{\frac{\kappa}{\sigma T}=\frac{\pi^2}{3}\Bigl(\frac{k_{\mathrm B}}{e}\Bigr)^2=\qty{2.44e-8}{W.\ohm/K^2}}{W·Ω/K$^2$ ＝ (J/K)$^2$ ÷ C$^2$}
\label{ph:原-8:WF}となります。これを\textbf{ウィーデマン・フランツの法則}といい、右辺の数を\textbf{ローレンツ数}といいます。\bkidx{うぃーでまんふらんつのほうそく}{ウィーデマン・フランツの法則}金属の種類によらない数で、抵抗率をはかれば熱伝導率がわかります（例題10）。

\begin{phyomi}[$\kappa/(\sigma T)=\tfrac{\pi^2}{3}(k_{\mathrm B}/e)^2$]
「熱の通しやすさは、電気の通しやすさと温度の積に比例し、比例の数は金属によらない」と読みます。電気は電荷 $e$ を、熱はエネルギー $k_{\mathrm B}T$ くらいを、同じ電子が同じ平均自由行程で運ぶので、比は $(k_{\mathrm B}/e)^2$ と温度だけで決まるのです。ドルーデの古典のモデルでは、係数が $\tfrac32$ になり、$\tfrac{\pi^2}{3}\approx3.3$ の半分以下でした。電子の比熱が小さすぎることと、速さが大きすぎることの2つのまちがいが、ほとんど打ち消し合っていたのです。
\end{phyomi}

\Needspace{14\baselineskip}\bksubsection{11}{くし形の結晶の帯を、プログラムで求める}
\label{ph:原-8:sec11}
⑦の帯とすき間を、数で確かめます。イオンが作るでこぼこを、間かく $a$ ごとに立った、細くて高い壁（くしの歯）で表します。この「くし形の結晶」は、クローニッヒとペニーが考えた、帯が手で計算できる、いちばん簡単なモデルです。
\begin{phmodel}[くし形の結晶（クローニッヒ・ペニー）]
\phignore{イオンのでこぼこの細かい形（くしの歯の強さ1つにまとめる）。電子どうしの反発。3次元の向き。}
\phassume{1次元で、間かく $a$ ごとに、同じ強さの細い壁が無限に並ぶ。歯と歯のあいだでは、電子は自由。}
\phlimit{本物の3次元の結晶（帯の形は向きによって違う）。原子の並びがくずれた物（ガラスのような非晶質や、表面）。}
\end{phmodel}

歯と歯のあいだでは、電子は自由な波 $\sin qx$・$\cos qx$ で、エネルギーは $E=\dfrac{\hbar^2q^2}{2m_{\mathrm e}}$ です。歯の所で $\psi$ と $\dd\psi/\dd x$ のつながり方の条件を入れ、さらに、並びが周期的なので、波は1つの区間ごとに同じ形をくり返し、ずれは「$e^{iKa}$ 倍」だけ（$i$ は虚数単位：原-5 ④。$K$ はこの結晶の中の波の番号）という条件を使います。計算すると（専門の本で扱います）、
\pheq{\cos Ka=\cos qa+P\,\frac{\sin qa}{qa}}{1 ＝ 1}
\label{ph:原-8:KP}が得られます。$P$ は歯の強さを表す数（単位のない数）で、$P=0$ なら歯がなく、自由電子です。

\begin{phyomi}[$\cos Ka=\cos qa+P\,\dfrac{\sin qa}{qa}$]
左辺は $\cos$ なので、$-1$ から $1$ の値しかとれません。右辺は、エネルギー（$q$）で決まる数です。右辺が $-1$ から $1$ のあいだに入るエネルギーでは、結晶の中を進む波があります（\textbf{帯}）。右辺が1より大きいか $-1$ より小さいエネルギーでは、そのような $K$ がなく、進む波がありません（\textbf{すき間}）。$qa=\pi,\ 2\pi,\ldots$ では $\sin qa=0$ で右辺が $\pm1$ になり、帯の上のはしになります。これは⑦のブラッグ反射の条件 $q=\pi/a$ と同じで、エネルギーは $h^2/(8m_{\mathrm e}a^2)$ の 1・4・9 倍です。
\end{phyomi}

\begin{bkfigure}
\begin{tikzpicture}
\begin{axis}[bkplot, axis lines=left, width=106mm, height=56mm, xmin=0, xmax=3.3, ymin=-2.3, ymax=6.2, xtick={1,2,3}, xticklabels={$\pi$,$2\pi$,$3\pi$}, ytick={-1,0,1}, xlabel={$qa$}, ylabel={右辺}]
  \fill[bkGreen!18] (axis cs:0.7247,-2.3) rectangle (axis cs:1,6.2);
  \fill[bkGreen!18] (axis cs:1.5004,-2.3) rectangle (axis cs:2,6.2);
  \fill[bkGreen!18] (axis cs:2.3577,-2.3) rectangle (axis cs:3,6.2);
  \draw[bkGray, line width=0.6pt, dashed] (axis cs:0,1) -- (axis cs:3.3,1) (axis cs:0,-1) -- (axis cs:3.3,-1);
  \addplot[bkcurve, domain=0.03:3.25, samples=300] {cos(deg(pi*x)) + 4.712389*sin(deg(pi*x))/(pi*x)};
  \node[font=\sffamily\footnotesize, anchor=south] at (axis cs:0.862,3.2) {帯1};
  \node[font=\sffamily\footnotesize, anchor=south] at (axis cs:1.75,3.2) {帯2};
  \node[font=\sffamily\footnotesize, anchor=south] at (axis cs:2.68,3.2) {帯3};
\end{axis}
\end{tikzpicture}
\bkfigcaption{くし形の結晶（$P=3\pi/2$）の右辺の値（実線）。横の $qa$ はエネルギーの平方根に比例する。右辺が $-1$ と1（破線）のあいだに入るエネルギー（色をぬった3つの帯）だけ、結晶の中を進む波がある。そのあいだが、すき間。帯の上のはしは、どれも $qa=\pi,2\pi,3\pi$。}\label{ph:原-8:fig-kp}
\end{bkfigure}

\phboxhead{プログラムで、帯のはしを求める}
$P=3\pi/2$、$a=\qty{0.5}{nm}$ として、$qa$ を細かく変えながら右辺を計算し、$|\text{右辺}|\le1$ となる範囲のはしを探します（例題8）。
\begin{phcode}[Python]
import numpy as np
P, a = 3 * np.pi / 2, 0.5e-9                  \# くしの強さ、くしの間かく
hbar, m, e = 1.0546e-34, 9.109e-31, 1.602e-19
qa = np.linspace(1e-6, 3.2 * np.pi, 200000)
F = np.cos(qa) + P * np.sin(qa) / qa          \# 右辺の関数
E = (hbar * qa / a)**2 / (2 * m) / e          \# その波のエネルギー（eV）
ok = np.abs(F) <= 1                           \# 帯の中（波が進める）
edge = np.flatnonzero(np.diff(ok.astype(int)))
print(np.round(E[edge], 2))                   \# 帯のはしのエネルギー
\end{phcode}
\begin{phcodeout}[実行結果]
[ 0.77  1.5   3.38  6.02  8.38 13.54]
\end{phcodeout}
帯は 0.77〜1.50 eV、3.38〜6.02 eV、8.38〜13.54 eV で、すき間は 1.88 eV と 2.36 eV です。帯の上のはし 1.50・6.02・13.54 eV は、自由電子の $h^2/(8m_{\mathrm e}a^2)$ の 1・4・9 倍と同じです。

\begin{phbaisuru}[大きく]{くしの強さ $P$}
$P$ を $\pi/2$・$3\pi/2$・$5\pi$ と大きくすると、いちばん下の帯は 0.38〜1.50 eV、0.77〜1.50 eV、1.19〜1.50 eV と、上のはしは変わらずに、はばがせまくなります。歯が強いほど、電子はとなりの区間へ移りにくく、それぞれの区間に閉じこめられた電子（原-5 の箱の電子：幅 0.5 nm で 1.50 eV）に近づくのです。反対に $P\to0$ では、すき間が消えて、自由電子にもどります。原子どうしの電子の雲の重なりが小さいほど、帯がせまくなる、という⑦の見方1と同じことです。
\end{phbaisuru}

\phboxhead{原子のゆれの比熱：デバイの $T^3$ 則}
熱-15 ⑧のアインシュタインのモデルでは、固体の原子がみな同じ角振動数 $\omega$ でゆれるとしたので、低い温度で比熱が $e^{-\theta_{\mathrm E}/T}$ のように急に0になりました。本物の固体の比熱は、もっとゆっくり、$T^3$ に比例して0に近づきます。デバイは、原子のゆれを、結晶全体を伝わる音の波（格子の波）の集まりと考えました。波の席は、④の電子と同じように、波数の空間の球で数えられます。低い温度で熱のエネルギーを受け取れるのは、$\hbar\omega\lesssim k_{\mathrm B}T$ の、長い波長（小さな波数）の波だけで、その数は球の体積に比例して $T^3$ に比例します。1つの波が受け取るエネルギーが $k_{\mathrm B}T$ くらいなので、エネルギーは $T^4$、比熱は $T^3$ に比例します。くわしく計算すると、低い温度で
\pheq{C=\frac{12\pi^4}{5}\,R\,\Bigl(\frac{T}{\theta_{\mathrm D}}\Bigr)^3}{J/(mol·K) ＝ J/(mol·K)}
\label{ph:原-8:debye}です。$\theta_{\mathrm D}$ は\textbf{デバイ温度}で、物によって決まる温度です（銅で約 343 K）。\bkidx{でばいおんど}{デバイ温度}この式は $T$ が $\theta_{\mathrm D}$ の10分の1ほどより低いときに使えます。高い温度では、比熱は $3R$（デュロン・プティの値）に近づきます。

\begin{bkfigure}
\begin{tikzpicture}
\begin{axis}[bkplot, axis lines=left, width=112mm, height=52mm, xmin=0, xmax=310, ymin=0, ymax=28, xtick={0,50,100,150,200,250,300}, ytick={0,5,10,15,20,24.94}, yticklabels={0,5,10,15,20,$3R$}, xlabel={温度（K）}, ylabel={$C$（J/(mol·K)）}]
  \draw[bkBlue, line width=1.7pt] (axis cs:2,0.000) -- (axis cs:6,0.010) -- (axis cs:10,0.048) -- (axis cs:14,0.132) -- (axis cs:18,0.281) -- (axis cs:22,0.513) -- (axis cs:26,0.844) -- (axis cs:30,1.287) -- (axis cs:34,1.845) -- (axis cs:38,2.511) -- (axis cs:40,2.880) -- (axis cs:50,4.985) -- (axis cs:60,7.276) -- (axis cs:70,9.495) -- (axis cs:80,11.506) -- (axis cs:90,13.263) -- (axis cs:100,14.766) -- (axis cs:110,16.042) -- (axis cs:120,17.120) -- (axis cs:130,18.034) -- (axis cs:140,18.809) -- (axis cs:150,19.469) -- (axis cs:160,20.035) -- (axis cs:170,20.521) -- (axis cs:180,20.942) -- (axis cs:190,21.308) -- (axis cs:200,21.628) -- (axis cs:210,21.909) -- (axis cs:220,22.157) -- (axis cs:230,22.376) -- (axis cs:240,22.571) -- (axis cs:250,22.745) -- (axis cs:260,22.901) -- (axis cs:270,23.041) -- (axis cs:280,23.168) -- (axis cs:290,23.282) -- (axis cs:300,23.386);
  \draw[bkOrange, line width=1.7pt, dashed] (axis cs:4,0.000) -- (axis cs:8,0.000) -- (axis cs:12,0.000) -- (axis cs:16,0.000) -- (axis cs:20,0.005) -- (axis cs:24,0.033) -- (axis cs:28,0.125) -- (axis cs:32,0.330) -- (axis cs:36,0.681) -- (axis cs:40,1.189) -- (axis cs:50,3.051) -- (axis cs:60,5.390) -- (axis cs:70,7.793) -- (axis cs:80,10.025) -- (axis cs:90,11.992) -- (axis cs:100,13.680) -- (axis cs:110,15.112) -- (axis cs:120,16.320) -- (axis cs:130,17.340) -- (axis cs:140,18.203) -- (axis cs:150,18.938) -- (axis cs:160,19.565) -- (axis cs:170,20.103) -- (axis cs:180,20.568) -- (axis cs:190,20.972) -- (axis cs:200,21.324) -- (axis cs:210,21.633) -- (axis cs:220,21.905) -- (axis cs:230,22.145) -- (axis cs:240,22.359) -- (axis cs:250,22.550) -- (axis cs:260,22.720) -- (axis cs:270,22.874) -- (axis cs:280,23.012) -- (axis cs:290,23.137) -- (axis cs:300,23.251);
  \addplot[bkGreen!60!black, line width=1.1pt, densely dotted, domain=0:80, samples=60] {1943.77*(x/343)^3};
  \draw[bkGray, line width=0.4pt] (axis cs:0,24.94) -- (axis cs:310,24.94);
  \node[font=\sffamily\footnotesize, anchor=west] at (axis cs:150,15.8) {デバイ（実線）};
  \node[font=\sffamily\footnotesize, anchor=west] at (axis cs:100,9.2) {アインシュタイン（破線）};
  \node[font=\sffamily\footnotesize, anchor=west] at (axis cs:80,22.5) {$T^3$ 則（点線）};
\end{axis}
\end{tikzpicture}
\bkfigcaption{銅の原子のゆれの比熱（$\theta_{\mathrm D}=\qty{343}{K}$ のデバイのモデル）。低い温度では $T^3$ 則（点線）に重なり、高い温度では $3R$ に近づく。アインシュタインのモデル（破線。$\theta_{\mathrm E}=0.806\,\theta_{\mathrm D}$ とした）は、低い温度で小さくなりすぎる。}\label{ph:原-8:fig-debye}
\end{bkfigure}

\begin{phbaisuru}{温度 $T$（低い温度で）}
$T^3$ に比例するので、$2^3=8$ 倍になります。銅の 10 K の比熱 \qty{0.048}{J/(mol.K)} は、20 K では \qty{0.39}{J/(mol.K)}（例題11）。逆に、温度を半分にすると8分の1です。とても低い温度では、原子のゆれの比熱がとても小さくなるので、⑥の電子の比熱（$T$ に比例）のほうが大きくなります。銅では、約3 K より下で入れかわります。
\end{phbaisuru}

\Needspace{12\baselineskip}\bksubsection{12}{帯の上に咲く、いろいろな性質}
\label{ph:原-8:sec12}
\phboxhead{超伝導}
ある種の金属や合金を、とても低い温度（\textbf{臨界温度} $T_{\mathrm c}$）より冷やすと、電気抵抗がまったく0になります。これを\textbf{超伝導}といいます。\bkidx{ちょうでんどう}{超伝導}1911年にオネスが水銀で見つけました。しくみは、1957年にバーディーン・クーパー・シュリーファーが説明しました（BCS 理論。1972年のノーベル物理学賞）。$E_{\mathrm F}$ の近くの電子が、結晶の原子のゆれを仲立ちにして、スピンを反対にした2個ずつの組（クーパー対）を作り、たくさんの組が1つの量子の状態にそろって動きます。組をこわすには、$E_{\mathrm F}$ の周りにできたすき間 $2\Delta$ のエネルギーが要り、BCS 理論では
\pheq{\Delta=1.76\,k_{\mathrm B}T_{\mathrm c}}{J ＝ J}
\label{ph:原-8:bcs}です。ふつうの散らばりでは、このすき間をこえるエネルギーを電子に与えられないので、電流が減らずに流れ続けます。病院の MRI の強い磁場（1.5〜3 T：電-6）は、ニオブとチタンの合金（Nb-Ti。$T_{\mathrm c}$ は約9 K）の線を液体ヘリウム（4.2 K）で冷やした、超伝導のコイルでつくっています（例題12）。

\begin{column}{安全：超伝導の実験は、見るだけ}
超伝導の磁石が浮く実験では、液体窒素（$-196$ ℃。窒素の沸点は 77 K）でセラミックスの超伝導体を冷やします。液体窒素は、皮ふにつくと、ひどいしもやけ（凍傷）になります。また、液体窒素が気体になると体積が何百倍にもふくらみ、せまい部屋では空気中の酸素がうすくなって危険です。ふたをしめた容器に入れると、破裂することもあります。液体窒素を使う実験は、学校の先生や科学館の人が、決まった方法で行うのを見るだけにしましょう。
\end{column}

\phboxhead{強誘電体}
チタン酸バリウム（BaTiO$_3$）の結晶では、酸素の八面体の中のチタンのイオンが、まん中から少しずれた所で落ち着きます。すべての小さな単位で同じ向きにずれるので、電場をかけなくても電気の偏り（分極）が残り、その向きを電場で反対にできます。これが電-3 の\textbf{強誘電体}で、ずれる向きは、となりの単位との電気の力と、電子の結合の席（原-7 ⑪）のかね合いで決まります。積層セラミックコンデンサーや、圧電素子（超音波・着火装置）に使われています。

\phboxhead{熱電効果}
金属の線の片方を温めると、熱いほうの $E_{\mathrm F}$ の近くの電子のフェルミ分布のすそが広がり（⑤）、電子が冷たいほうへ流れ出そうとして、電圧ができます（電-4 の熱電対）。電子の比熱と同じく、関われるのは $E_{\mathrm F}$ の近くの電子だけなので、金属の電圧は $(k_{\mathrm B}/e)\times(k_{\mathrm B}T/E_{\mathrm F})$ くらい、1 K あたり数 \si{\micro V}と小さくなります。K 熱電対の 41 \si{\micro V}/K（電-4）は、合金の $d$ の帯の形の違いから生まれる差です。半導体では電子の数がずっと少なく、1 K あたり数百 \si{\micro V}にもなります（ペルチェ素子：原-9 ⑬）。

\phboxhead{磁石：帯の中の強磁性}
鉄の $3d$ の電子は、金属の中でも、となりの原子と席を分け合って帯になります。原-7 ⑧の交換相互作用のために、上向きのスピンの帯と下向きのスピンの帯が、少しずれたエネルギーに分かれ、上向きの電子のほうが多くなります。その差が、原子1個あたり約 $2.2\mu_{\mathrm B}$ の、整数でない磁気モーメントです（原-7 の例題9）。これが、電-7 の磁区の中で向きがそろう強磁性のしくみです。

\begin{senmonsho}[固体物理の専門書の書き方]
\begin{bkhonbun}
周期ポテンシャル $V(\symbfit r+\symbfit R)=V(\symbfit r)$ 中の電子の固有関数は，ブロッホの定理により $\psi_{n\symbfit k}(\symbfit r)=e^{i\symbfit k\cdot\symbfit r}u_{n\symbfit k}(\symbfit r)$（$u_{n\symbfit k}$ は格子の周期をもつ）と書ける．エネルギー固有値 $E_n(\symbfit k)$ をバンド構造とよび，$\symbfit k$ は第1ブリルアン帯域内にとる．ブリルアン帯域の境界でブラッグ反射が起こり，エネルギーギャップが開く．
\end{bkhonbun}
\bkyaku{「周期ポテンシャル」は、この章のイオンのでこぼこ。$\symbfit R$ は結晶の並びをずらす矢印（太字はベクトル）。ブロッホの定理は、⑪の「区間ごとに $e^{iKa}$ 倍」の3次元版で、$n$ は帯の番号、$\symbfit k$ はこの章の $K$。ブリルアン帯域は、1次元なら $-\pi/a\le K\le\pi/a$ の範囲。「エネルギーギャップ」はすき間（禁制帯）。}
\end{senmonsho}

\begin{phmada}[高温超伝導のしくみは？]
1986年、ベドノルツとミュラーは、銅の酸化物のセラミックスが、それまでの超伝導体より高い温度で超伝導になることを見つけました（1987年のノーベル物理学賞）。その後、液体窒素の温度（77 K）より高い $T_{\mathrm c}$ の物質も見つかりました。けれども、なぜこんなに高い温度で超伝導になるのか、そのしくみは、⑫の BCS 理論のような「原子のゆれを仲立ちにした組」では説明しきれていません。電子どうしの強い反発が主役らしい、と考えられていますが、何千もの論文が出た今も、研究者の意見は一つにまとまっていません（「深い未解決の問題」と、2015年の総説に書かれています）。もし室温で超伝導になる物質を作れれば、電気を損失なく送れるようになると期待されています。
\end{phmada}

\begin{phsaki}{帯の物理から、デバイスと材料へ}
この章の帯とすき間は、固体物理学（物性物理学）の出発点です。3次元の本物の結晶の帯は、コンピュータで計算して（第一原理計算）、向きごとの $E$ と $K$ の関係（バンド構造）として描きます。半導体の材料の設計、太陽電池・LED・パワー半導体（原-9・原-10）、超伝導の線材、磁気記録の材料、熱電変換材料などは、すべてこの土台の上で作られています。くわしくは、固体物理・半導体工学・材料工学・超伝導工学の専門の本で学びます。
\end{phsaki}



\begin{taisetsu}
\def\tsl#1{\makebox[7.2\zw][l]{\sffamily\bfseries\small\color{\bkTT}#1}}
\tsl{ブラッグ}$2d\sin\theta=n\lambda$\\[3pt]
\tsl{フェルミ}$E_{\mathrm F}=\dfrac{\hbar^2}{2m_{\mathrm e}}(3\pi^2n)^{2/3}$、$f(E)=\dfrac{1}{e^{(E-E_{\mathrm F})/(k_{\mathrm B}T)}+1}$\\[3pt]
\tsl{すき間}$E=\dfrac{h^2}{8m_{\mathrm e}a^2}$（$\lambda=2a$）、上がる割合 $e^{-E_{\mathrm g}/(2k_{\mathrm B}T)}$\\[3pt]
\tsl{電子の比熱}$C_{\text{電子}}=\dfrac{\pi^2}{2}Nk_{\mathrm B}\dfrac{T}{T_{\mathrm F}}$\\[3pt]
\tsl{運ぶ性質}$\ell=v_{\mathrm F}\tau$、$\dfrac{\kappa}{\sigma T}=\dfrac{\pi^2}{3}\Bigl(\dfrac{k_{\mathrm B}}{e}\Bigr)^2$\\[3pt]
\tsl{格子・超伝導}$C=\dfrac{12\pi^4}{5}R\Bigl(\dfrac{T}{\theta_{\mathrm D}}\Bigr)^3$、$\Delta=1.76\,k_{\mathrm B}T_{\mathrm c}$
\end{taisetsu}

\begin{kotoba}
\small
\bkword{結晶}{原子が、同じ形を繰り返して規則正しく並んだ固体。}
\bkword{エネルギーバンド（帯）}{たくさんの原子の席が集まってできた、ぎっしりつまったエネルギーの範囲。}
\bkword{禁制帯（すき間）}{帯と帯のあいだの、電子の席のないエネルギーの範囲。そのはばが禁制帯幅 $E_{\mathrm g}$。}
\bkword{フェルミエネルギー}{0 K で、電子がうめているいちばん上の席のエネルギー。}
\bkword{価電子帯・伝導帯}{半導体・絶縁体で、電子でうまった下の帯と、空いた上の帯。}
\bkword{平均自由行程}{電子が、並びの乱れにぶつかるまでに進む平均の距離。}
\bkword{超伝導}{臨界温度より低い温度で、電気抵抗が0になること。}
\end{kotoba}

\bkkigocards{\bkkigocard{$E_{\mathrm F}$}{イー・エフ}{フェルミエネルギー〔eV〕。熱-15 の化学ポテンシャル $\mu$（0 K）と同じもの}{Fermi energy}{原-8}}
{\bkkigocard{$E_{\mathrm g}$}{イー・ジー}{禁制帯幅〔eV〕。価電子帯の上と伝導帯の底の差}{band gap}{原-8}}
\bkkigocards{\bkkigocard{$\theta_{\mathrm D}$}{シータ・ディー}{デバイ温度〔K〕。熱-15 のアインシュタイン温度 $\theta_{\mathrm E}$ とは別の量}{Debye temperature}{原-8}}
{\bkkigocard{$\kappa$}{カッパ}{熱伝導率〔W/(m·K)〕。熱-11 の $k$ と同じ量。原-6 の減衰の定数 $\kappa$ とは別}{thermal conductivity}{原-8}}

%% ======================== 例題 ========================
\Needspace{10\baselineskip}\bksection{reidai}{例題でたしかめよう}
\label{ph:原-8:ex1}
\begin{phreidai}[ケイ素とガラス（はてな？）]
禁制帯幅は、ケイ素が \qty{1.12}{eV}、ガラス（二酸化ケイ素）が約 \qty{9}{eV}です。300 K（$k_{\mathrm B}T=\qty{0.0259}{eV}$）で、上の帯へ上がれる電子の割合の目安 $e^{-E_{\mathrm g}/(2k_{\mathrm B}T)}$ を、それぞれ求めて比べましょう。
\phbamen
半導体と絶縁体の、電気の通しやすさの違いのもとを、数で比べる場面です。
\phmushi
上の帯・下の帯の席の数の違い（式の前にかかる係数）は無視し、指数の部分だけで比べます。
\phzu
$E_{\mathrm g}=\qty{1.12}{eV}$ と \qty{9.0}{eV}、$k_{\mathrm B}T=\qty{0.0259}{eV}$。
\phriyu
$E_{\mathrm F}$ はすき間のまん中にあり、上の帯の底までは $E_{\mathrm g}/2$。そこの席に電子がいる確率は、フェルミ分布のすそ（⑤・⑧）です。
\phtateru
\[ f=e^{-E_{\mathrm g}/(2k_{\mathrm B}T)} \]
\phtoku
\[ \text{ケイ素：}e^{-1.12/(2\times0.0259)}=e^{-21.6}\approx4.1\times10^{-10} \]
\[ \text{ガラス：}e^{-9.0/(2\times0.0259)}=e^{-173.7}\approx3\times10^{-76} \]
比は約 $10^{66}$ です（$\log_{10}$ で $66.1$）。
\phtashikame
単位：指数の中は eV÷eV で、単位のない数。\\
大きさ：ケイ素でも100億分の4ほどしか上がりませんが、1 m$^3$ に $10^{29}$ 個ほどある価電子のうちの話なので、上がる電子の数は何十億個にもなります。ガラスでは1個もいません。\\
極端な場合：$E_{\mathrm g}\to0$ なら割合は1で、金属と同じようになります。
\kotae ケイ素 約 $4.1\times10^{-10}$、ガラス 約 $3\times10^{-76}$。ケイ素のほうが約 $10^{66}$ 倍大きい
\end{phreidai}
\begin{jibun}
LED や太陽電池に使われるヒ化ガリウム（GaAs）の禁制帯幅を \qty{1.42}{eV}とします。300 K での $e^{-E_{\mathrm g}/(2k_{\mathrm B}T)}$ を求め、ケイ素の何倍かを答えましょう。
\bkjans{$e^{-1.42/0.0518}=e^{-27.4}\approx1.2\times10^{-12}$。ケイ素の約0.003倍（約300分の1）}
\end{jibun}

\label{ph:原-8:ex2}
\begin{phreidai}[ケイ素の結晶の X線回折（3けたで書く）]
ケイ素の結晶の、立方体の対角線に垂直な面（間かく $d=a/\sqrt3$、$a=\qty{0.5431}{nm}$）に、銅の Kα$_1$ 線（波長 \qty{0.1541}{nm}）を当てます。$n=1$ で強く反射する角 $\theta$ と、装置で読む角 $2\theta$（入ってくる向きと出ていく向きの角）を求めましょう（装置の目もりと比べるので、3けたで書きます）。
\phbamen
X線の回折の装置で、ケイ素の結晶の面を調べる場面です。
\phmushi
X線は1つの波長だけとし、結晶は完全に規則正しいとします。
\phzu
$d=0.5431/\sqrt3\approx\qty{0.3136}{nm}$、$\lambda=\qty{0.1541}{nm}$、$n=1$。
\phriyu
平行な原子の面での反射なので、ブラッグの条件（②）を使います。
\phtateru
\[ \sin\theta=\frac{n\lambda}{2d} \]
\phtoku
\[ \sin\theta=\frac{0.1541}{2\times0.3136}\approx0.2457,\qquad \theta\approx14.2^\circ,\qquad 2\theta\approx28.4^\circ \]
\phtashikame
単位：nm÷nm で、$\sin$ は単位のない数。\\
大きさ：ケイ素の粉の X線回折で、いちばん強く出る線の 28.4° と合います。\\
極端な場合：$n=2$ なら $\sin\theta=0.491$、$\theta\approx29.4^\circ$。$n\lambda>2d$ になる $n$（$n=5$ 以上）では、反射は起きません。
\kotae $\theta\approx14.2^\circ$、$2\theta\approx28.4^\circ$
\end{phreidai}
\begin{jibun}
ケイ素の別の面（間かく $d=a/\sqrt8$）では、銅の Kα$_1$ 線は $2\theta$ が何度の所で強く反射しますか（3けたで書く）。
\bkjans{$d\approx\qty{0.1920}{nm}$。$\sin\theta=0.1541/0.3840\approx0.4013$、$\theta\approx23.7^\circ$、$2\theta\approx47.3^\circ$}
\end{jibun}

\label{ph:原-8:ex3}
\begin{phreidai}[ケイ素の格子定数から、密度と原子の数を求める]
ケイ素（モル質量 \qty{28.09}{g/mol}）はダイヤモンド形で、格子定数は $a=\qty{0.5431}{nm}$ です。密度（g/cm$^3$）と、1 m$^3$ あたりの原子の数を求めましょう。$N_{\mathrm A}=\num{6.022e23}\,\mathrm{/mol}$ とします。
\phbamen
X線ではかった原子の並びから、手で持てる大きさの物の性質を求める場面です。
\phmushi
結晶には欠陥や不純物がなく、どこでも同じ並びだとします。
\phzu
立方体1つに8個（①）。$a^3=(0.5431\times10^{-9})^3\approx\qty{1.602e-28}{m^3}$。
\phriyu
密度は「質量÷体積」。立方体1つの質量は、原子8個の質量です（③）。
\phtateru
\[ \rho=\frac{8\times M/N_{\mathrm A}}{a^3},\qquad n=\frac{8}{a^3} \]
\phtoku
\[ \rho=\frac{8\times28.09\times10^{-3}\div(6.022\times10^{23})}{1.602\times10^{-28}}\approx\qty{2.33e3}{kg/m^3}=\qty{2.33}{g/cm^3} \]
\[ n=\frac{8}{1.602\times10^{-28}}\approx5.0\times10^{28}\ \text{個/m}^3 \]
\phtashikame
単位：kg÷m$^3$。\\
大きさ：天びんではかったケイ素の密度 \qty{2.330}{g/cm^3}と合います。原子の数は銅（$8.5\times10^{28}$）より少し少なく、すき間の多いダイヤモンド形の並びだからです（原子どうしは近くても、つまり方がゆるい）。\\
極端な場合：同じ $a$ で面心立方（4個）なら、密度も数も半分です。
\kotae 密度 約 \qty{2.33}{g/cm^3}、原子の数 約 $5.0\times10^{28}$ 個/m$^3$
\end{phreidai}
\begin{jibun}
銅は面心立方（立方体1つに4個）で、密度は \qty{8.96}{g/cm^3}、モル質量は \qty{63.55}{g/mol}です。格子定数 $a$ を求めましょう。
\bkjans{$a^3=4M/(\rho N_{\mathrm A})\approx\qty{4.71e-29}{m^3}$。$a\approx\qty{0.361}{nm}$}
\end{jibun}

\label{ph:原-8:ex4}
\begin{phreidai}[銅のフェルミエネルギーと、電子の速さ]
銅の自由電子の数は $n=8.49\times10^{28}$ 個/m$^3$ です（電-4）。フェルミエネルギー（eV）、フェルミ温度、フェルミ速度を求めましょう。$\hbar=\qty{1.055e-34}{J.s}$、$m_{\mathrm e}=\qty{9.11e-31}{kg}$、$k_{\mathrm B}=\qty{1.38e-23}{J/K}$ とします。
\phbamen
電-4 で「1秒に千 km 以上」と書いた電子の速さを求める場面です。
\phmushi
イオンの並びのでこぼこを無視し、自由電子とします（④のモデル化の札）。
\phzu
$n=8.49\times10^{28}$ /m$^3$。
\phriyu
$N$ 個の電子を3次元の箱の席に低い順につめたときの、いちばん上の席のエネルギーが $E_{\mathrm F}$（④）。速さは $\tfrac12m_{\mathrm e}v_{\mathrm F}^{\,2}=E_{\mathrm F}$ から。
\phtateru
\[ k_{\mathrm F}=(3\pi^2n)^{1/3},\qquad E_{\mathrm F}=\frac{\hbar^2k_{\mathrm F}^{\,2}}{2m_{\mathrm e}},\qquad T_{\mathrm F}=\frac{E_{\mathrm F}}{k_{\mathrm B}},\qquad v_{\mathrm F}=\frac{\hbar k_{\mathrm F}}{m_{\mathrm e}} \]
\phtoku
\[ k_{\mathrm F}=(3\pi^2\times8.49\times10^{28})^{1/3}\approx\num{1.36e10}\,\mathrm{/m} \]
\[ E_{\mathrm F}=\frac{(1.055\times10^{-34}\times1.36\times10^{10})^2}{2\times9.11\times10^{-31}}\approx\qty{1.13e-18}{J}\approx\qty{7.0}{eV} \]
\[ T_{\mathrm F}=\frac{1.13\times10^{-18}}{1.38\times10^{-23}}\approx\qty{8.2e4}{K} \]
\[ v_{\mathrm F}=\frac{1.055\times10^{-34}\times1.36\times10^{10}}{9.11\times10^{-31}}\approx\qty{1.57e6}{m/s} \]
\phtashikame
単位：(J·s/m)$^2$÷kg＝J、(J·s/m)÷kg＝m/s。\\
大きさ：1秒に約 \qty{1570}{km}。光の速さの約0.5\,\% なので、相対論（原-11）は要りません。300 K の気体の分子のように考えた速さ（$\sqrt{3k_{\mathrm B}T/m_{\mathrm e}}\approx\qty{117}{km/s}$）の13倍です。\\
極端な場合：温度によらず、0 K でもこの速さで動き回っています。
\kotae $E_{\mathrm F}\approx\qty{7.0}{eV}$、$T_{\mathrm F}\approx\qty{8.2e4}{K}$、$v_{\mathrm F}\approx\qty{1.57e6}{m/s}$（1秒に約 \qty{1570}{km}）
\end{phreidai}
\begin{jibun}
銀の自由電子の数は $5.86\times10^{28}$ 個/m$^3$ です（原-7 のじぶんでとこう8）。フェルミエネルギーとフェルミ速度を求めましょう。
\bkjans{$k_{\mathrm F}\approx\num{1.20e10}\,\mathrm{/m}$。$E_{\mathrm F}\approx\qty{5.5}{eV}$、$v_{\mathrm F}\approx\qty{1.39e6}{m/s}$}
\end{jibun}

\label{ph:原-8:ex5}
\begin{phreidai}[フェルミ分布で、席がうまる割合]
300 K（$k_{\mathrm B}T=\qty{0.0259}{eV}$）で、$E_{\mathrm F}$ より 0.05 eV 下・上、0.10 eV 下・上の席に、電子がいる確率を求めましょう。
\phbamen
$E_{\mathrm F}$ の周りで、段がどれくらいなだらかになっているかを、数で見る場面です。
\phmushi
$E_{\mathrm F}$ は温度でほとんど変わらないとします。
\phzu
$E-E_{\mathrm F}=-0.10,\ -0.05,\ +0.05,\ +0.10$ eV。
\phriyu
席に電子がいる確率は、フェルミ分布（⑤）。
\phtateru
\[ f=\frac{1}{e^{(E-E_{\mathrm F})/(k_{\mathrm B}T)}+1} \]
\phtoku
\[ E-E_{\mathrm F}=+0.05\ \text{eV}\ \text{のとき}\quad f=\frac{1}{e^{1.93}+1}=\frac{1}{6.89+1}\approx0.13 \]
\[ E-E_{\mathrm F}=-0.05\ \text{eV}\ \text{のとき}\quad f=\frac{1}{e^{-1.93}+1}=\frac{1}{0.145+1}\approx0.87 \]
\[ E-E_{\mathrm F}=+0.10\ \text{eV}\ \text{のとき}\quad f=\frac{1}{e^{3.86}+1}=\frac{1}{47.5+1}\approx0.021 \]
\[ E-E_{\mathrm F}=-0.10\ \text{eV}\ \text{のとき}\quad f=\frac{1}{e^{-3.86}+1}\approx0.98 \]
\phtashikame
単位：指数の中は単位のない数。\\
大きさ：$E_{\mathrm F}$ から上下に $0.1$ eV（$k_{\mathrm B}T$ の4倍）はなれると、ほぼ1か0です。\\
極端な場合：同じだけ上と下の値をたすと、$0.13+0.87=1$、$0.021+0.98\approx1$ です。$E_{\mathrm F}$ の上に上がった電子の数と、下に空いた席の数が等しいことを表しています。
\kotae 0.05 eV 下で 0.87、上で 0.13。0.10 eV 下で 0.98、上で 0.021
\end{phreidai}
\begin{jibun}
1000 K（$k_{\mathrm B}T=\qty{0.0862}{eV}$）で、$E_{\mathrm F}$ より 0.10 eV 上の席に電子がいる確率を求めましょう。
\bkjans{$1/(e^{1.16}+1)\approx0.24$}
\end{jibun}

\label{ph:原-8:ex6}
\begin{phreidai}[銅の電子の比熱は、等分配の何分の1？]
銅のフェルミ温度は $T_{\mathrm F}=\qty{8.2e4}{K}$ です（例題4）。300 K での電子の比熱は、等分配の法則で考えた値 $\tfrac32Nk_{\mathrm B}$ の何倍ですか。
\phbamen
金属の比熱に電子の分が見えない理由を、数で確かめる場面です（⑥）。
\phmushi
自由電子のモデルを使い、原子のゆれの比熱とは別に考えます。
\phzu
$T=\qty{300}{K}$、$T_{\mathrm F}=\qty{8.2e4}{K}$。
\phriyu
熱を受け取れるのは $E_{\mathrm F}$ の近くの電子だけで、その比熱はゾンマーフェルトの式（⑥）です。
\phtateru
\[ \frac{C_{\text{電子}}}{\tfrac32Nk_{\mathrm B}}=\frac{\tfrac{\pi^2}{2}Nk_{\mathrm B}T/T_{\mathrm F}}{\tfrac32Nk_{\mathrm B}}=\frac{\pi^2}{3}\cdot\frac{T}{T_{\mathrm F}} \]
\phtoku
\[ \frac{\pi^2}{3}\times\frac{300}{8.2\times10^4}\approx3.29\times0.00366\approx0.012 \]
\phtashikame
単位：比なので単位はない。\\
大きさ：約80分の1。原子のゆれの比熱 $3R$ にくらべると、さらにその半分で、0.6\,\% ほどにしかなりません。\\
極端な場合：$T$ が $T_{\mathrm F}$ に近い（何万 K）ほどの高温なら、電子も等分配に近づきますが、そのころには銅はとけて蒸発しています。
\kotae 約 0.012 倍
\end{phreidai}
\begin{jibun}
30 K では、銅の電子の比熱は、等分配の値の何倍ですか。
\bkjans{$\frac{\pi^2}{3}\times\frac{30}{8.2\times10^4}\approx0.0012$ 倍}
\end{jibun}

\label{ph:原-8:ex7}
\begin{phreidai}[1次元の結晶で、すき間ができるエネルギー]
イオンが間かく $a=\qty{0.25}{nm}$ で1列に並んだ1次元の結晶の中で、自由電子の波がはじめてブラッグ反射して、すき間ができるのは、何 eV の所ですか。
\phbamen
⑦のすき間が、どれくらいのエネルギーにできるかを見積もる場面です。
\phmushi
すき間のはばは考えず、すき間の位置（まん中のエネルギー）だけを、自由電子の式で求めます。
\phzu
$\lambda=2a=\qty{0.50}{nm}$。
\phriyu
となりのイオンで反射した波がそろう条件は $2a=\lambda$（⑦）。ド・ブロイの関係 $p=h/\lambda$（原-3）と、$E=p^2/(2m_{\mathrm e})$ を使います。
\phtateru
\[ E=\frac{p^2}{2m_{\mathrm e}}=\frac{1}{2m_{\mathrm e}}\Bigl(\frac{h}{2a}\Bigr)^2=\frac{h^2}{8m_{\mathrm e}a^2} \]
\phtoku
\[ E=\frac{(6.63\times10^{-34})^2}{8\times9.11\times10^{-31}\times(0.25\times10^{-9})^2}\approx\qty{9.6e-19}{J}\approx\qty{6.0}{eV} \]
\phtashikame
単位：(J·s)$^2$÷(kg·m$^2$)＝J。\\
大きさ：銅の $E_{\mathrm F}$（7.0 eV）と同じくらい。金属の電子の波長は、原子の間かくの2倍くらいなので、帯の形は自由電子から少しずれます。\\
極端な場合：間かく $a$ を2倍にすると、エネルギーは4分の1。
\kotae 約 \qty{6.0}{eV}
\end{phreidai}
\begin{jibun}
間かくが $a=\qty{0.30}{nm}$ の1次元の結晶では、はじめのすき間は何 eV の所にできますか。
\bkjans{$h^2/(8m_{\mathrm e}a^2)\approx\qty{4.2}{eV}$}
\end{jibun}

\label{ph:原-8:ex8}
\begin{phreidai}[くし形の結晶の帯を、プログラムで求める]
くし形の結晶（$P=3\pi/2$、$a=\qty{0.5}{nm}$）で、⑪のプログラムを動かし、エネルギーの低いほうから3つの帯の範囲と、2つのすき間のはばを求めましょう。
\phbamen
帯とすき間が本当にできることを、手で解けない式から、数で確かめる場面です。
\phmushi
1次元のくし形の結晶のモデル（⑪のモデル化の札）。
\phzu
右辺 $F=\cos qa+P\dfrac{\sin qa}{qa}$、$E=\dfrac{\hbar^2q^2}{2m_{\mathrm e}}$。
\phriyu
$|F|\le1$ のエネルギーだけが帯（⑪）。$qa$ を細かく変えて、$|F|=1$ をまたぐ所を探します。
\phtateru
$qa$ を 0 から $3.2\pi$ まで20万に分けて、$|F|\le1$ か調べ、変わり目の $E$ を書き出す。
\phtoku
プログラムの結果：帯は 0.77〜1.50 eV、3.38〜6.02 eV、8.38〜13.54 eV。すき間は $3.38-1.50=\qty{1.88}{eV}$、$8.38-6.02=\qty{2.36}{eV}$。
\phtashikame
単位：eV。\\
大きさ：帯の上のはし 1.50・6.02・13.54 eV は、自由電子の $\tfrac{h^2}{8m_{\mathrm e}a^2}$（$a=\qty{0.5}{nm}$ で 1.50 eV）の1・4・9倍で、⑦の考えと合います。すき間は数 eV で、本物の結晶の禁制帯と同じくらいです。\\
極端な場合：$P=0$ なら、$F=\cos qa$ で、いつでも $|F|\le1$。すき間は消えます。
\kotae 帯 0.77〜1.50・3.38〜6.02・8.38〜13.54 eV、すき間 1.88 eV と 2.36 eV
\end{phreidai}
\begin{jibun}
例題8 のプログラムの1行めを \texttt{P = 5 * np.pi} にすると、いちばん下の帯は何 eV から何 eV になりますか。$P=3\pi/2$ のときと比べて、どうなりましたか。
\bkjans{1.19〜1.50 eV。上のはしは同じで、はばがせまくなった（電子が区間に閉じこめられやすくなった）}
\end{jibun}

\label{ph:原-8:ex9}
\begin{phreidai}[銅の電子の平均自由行程]
電-4 で求めた銅の $\tau=\qty{2.4e-14}{s}$ と、例題4 のフェルミ速度 $v_{\mathrm F}=\qty{1.57e6}{m/s}$ から、平均自由行程を求めましょう。銅の一番近い原子どうしの間かく（\qty{0.256}{nm}）の何倍ですか。
\phbamen
電子が、ぶつかるまでにどれだけ進むかを求める場面です（⑩）。
\phmushi
ぶつかり方の違いを、平均の時間 $\tau$ 1つにまとめます（電-4 のドルーデのモデル）。
\phzu
$\ell=v_{\mathrm F}\tau$。
\phriyu
電気を運ぶのは $E_{\mathrm F}$ の近くの電子で、速さは $v_{\mathrm F}$ です（⑥・⑩）。
\phtateru
\[ \ell=v_{\mathrm F}\tau,\qquad \frac{\ell}{0.256\ \text{nm}} \]
\phtoku
\[ \ell=1.57\times10^6\times2.4\times10^{-14}\approx\qty{3.8e-8}{m}=\qty{38}{nm},\qquad \frac{38}{0.256}\approx150 \]
\phtashikame
単位：m/s×s＝m。\\
大きさ：原子の約150個分。気体の分子のような速さ（117 km/s）で計算すると \qty{2.8}{nm}で、10個ほどしか進まない計算になり、「イオン1個1個にぶつかる」という古い考えと合ってしまいます。\\
極端な場合：とても純粋な銅を液体ヘリウムの温度まで冷やすと、原子のゆれが小さくなり、$\tau$ も $\ell$ も何百倍にもなります。
\kotae 約 \qty{38}{nm}、原子の間かくの約150倍
\end{phreidai}
\begin{jibun}
純粋な銅を冷やして、$\tau$ が100倍になったとすると、平均自由行程はいくらですか。
\bkjans{約 $\qty{3.8e-6}{m}$（\qty{3.8}{\micro m}）}
\end{jibun}

\label{ph:原-8:ex10}
\begin{phreidai}[抵抗率から、銅の熱伝導率を求める]
銅の 20 ℃（293 K）での抵抗率は $\rho=\qty{1.72e-8}{\ohm.m}$ です（電-4）。ウィーデマン・フランツの法則で、熱伝導率を求めましょう。
\phbamen
電気の性質から、熱の性質を求める場面です（⑩）。
\phmushi
熱を運ぶのは電子だけとし、原子のゆれ（格子の波）が運ぶ分は無視します。
\phzu
$\sigma=1/\rho$、$T=\qty{293}{K}$、ローレンツ数 $\qty{2.44e-8}{W.\ohm/K^2}$。
\phriyu
同じ電子が電気と熱を運ぶので、$\kappa=\dfrac{\pi^2}{3}\Bigl(\dfrac{k_{\mathrm B}}{e}\Bigr)^2\sigma T$（⑩）。
\phtateru
\[ \kappa=2.44\times10^{-8}\times\frac{T}{\rho} \]
\phtoku
\[ \kappa=2.44\times10^{-8}\times\frac{293}{1.72\times10^{-8}}\approx\qty{4.2e2}{W/(m.K)} \]
\phtashikame
単位：(W·\si{\ohm}/K$^2$)×K÷(\si{\ohm.m})＝W/(m·K)。\\
大きさ：熱-11 の表の銅の熱伝導率（約 400 W/(m·K)）と、5\,\% ほどの違いで合います。\\
極端な場合：抵抗率の大きいニクロム（\qty{1.1e-6}{\ohm.m}：電-4）では、熱伝導率は約 6.5 W/(m·K) で、銅の約60分の1です。
\kotae 約 \qty{4.2e2}{W/(m.K)}
\end{phreidai}
\begin{jibun}
アルミニウムの抵抗率は $\qty{2.65e-8}{\ohm.m}$ です（電-4）。293 K での熱伝導率を、ウィーデマン・フランツの法則で求めましょう。
\bkjans{約 $\qty{2.7e2}{W/(m.K)}$（熱-11 の表の 237 W/(m·K) より1割ほど大きい）}
\end{jibun}

\label{ph:原-8:ex11}
\begin{phreidai}[低い温度の銅の比熱：デバイの $T^3$ 則]
銅のデバイ温度を \qty{343}{K}とします。10 K と 20 K での、原子のゆれの比熱（1 mol あたり）を求めましょう。$R=\qty{8.31}{J/(mol.K)}$ とします。
\phbamen
とても低い温度での金属の比熱を見積もる場面です（⑪）。
\phmushi
電子の比熱は、原子のゆれの比熱と別に考え、ここでは入れません。
\phzu
$\theta_{\mathrm D}=\qty{343}{K}$、$T=10,\ \qty{20}{K}$。
\phriyu
$T$ が $\theta_{\mathrm D}$ の10分の1より低いので、デバイの $T^3$ 則（⑪）が使えます。
\phtateru
\[ C=\frac{12\pi^4}{5}R\Bigl(\frac{T}{\theta_{\mathrm D}}\Bigr)^3,\qquad \frac{12\pi^4}{5}R\approx\qty{1944}{J/(mol.K)} \]
\phtoku
\[ C_{10}=1944\times\Bigl(\frac{10}{343}\Bigr)^3\approx\qty{0.048}{J/(mol.K)} \]
\[ C_{20}=1944\times\Bigl(\frac{20}{343}\Bigr)^3\approx\qty{0.39}{J/(mol.K)} \]
\phtashikame
単位：J/(mol·K)。\\
大きさ：室温の $3R\approx\qty{25}{J/(mol.K)}$ の、500分の1と60分の1。\\
極端な場合：温度が2倍で比熱は8倍（$0.048\times8\approx0.39$）。
\kotae 10 K で 約 \qty{0.048}{J/(mol.K)}、20 K で 約 \qty{0.39}{J/(mol.K)}
\end{phreidai}
\begin{jibun}
例題11 の銅の、5 K での原子のゆれの比熱を求めましょう。
\bkjans{$1944\times(5/343)^3\approx\qty{0.0060}{J/(mol.K)}$}
\end{jibun}

\label{ph:原-8:ex12}
\begin{phreidai}[MRI の超伝導の線のすき間]
MRI の磁石の線に使う Nb-Ti の臨界温度を \qty{9.5}{K}とします。BCS 理論のすき間 $\Delta$ を meV の単位で求め、液体ヘリウムの温度（4.2 K）の $k_{\mathrm B}T$ と比べましょう。$k_{\mathrm B}=\qty{8.62e-5}{eV/K}$ とします。
\phbamen
超伝導の電子の組が、冷やした線の中で、熱でこわれないかを調べる場面です（⑫）。
\phmushi
BCS 理論の、いちばん簡単な場合の式を使います。
\phzu
$T_{\mathrm c}=\qty{9.5}{K}$、$T=\qty{4.2}{K}$。
\phriyu
すき間は $\Delta=1.76\,k_{\mathrm B}T_{\mathrm c}$（⑫）。
\phtateru
\[ \Delta=1.76\,k_{\mathrm B}T_{\mathrm c},\qquad k_{\mathrm B}\times4.2 \]
\phtoku
\[ \Delta=1.76\times8.62\times10^{-5}\times9.5\approx\qty{1.44e-3}{eV}\approx\qty{1.4}{meV} \]
\[ k_{\mathrm B}\times4.2=8.62\times10^{-5}\times4.2\approx\qty{3.6e-4}{eV}=\qty{0.36}{meV} \]
\phtashikame
単位：(eV/K)×K＝eV。\\
大きさ：すき間は、4.2 K の熱のエネルギーの約4倍。ボルツマン因子でおよそ $e^{-4}\approx0.02$ なので、組をこわせる熱の動きは少なく、電流は減らずに流れます。\\
極端な場合：温度を $T_{\mathrm c}$ に近づけると、すき間は0に向かって小さくなり（この式は 0 K の値）、超伝導が消えます。
\kotae $\Delta\approx\qty{1.4}{meV}$。4.2 K の $k_{\mathrm B}T$（約0.36 meV）の約4倍
\end{phreidai}
\begin{jibun}
臨界温度が \qty{4.0}{K}の超伝導体のすき間 $\Delta$ を、meV の単位で求めましょう。
\bkjans{$\Delta=1.76\,k_{\mathrm B}T_{\mathrm c}$ で、約 \qty{6.1e-4}{eV}（約 0.61 meV）}
\end{jibun}

%% ======================== まちがえやすいところ ========================
\Needspace{11\baselineskip}\bksection{tsumazuki}{まちがえやすいところをチェック}
\begin{tsumazuki}{電子は、イオンにぶつかりながら進む？}
\bkng{金属の電子は、となりのイオンに次々ぶつかるので、電気抵抗がある。}
\bkok{完全な結晶の中では、電子の波はイオンにぶつからない。ぶつかるのは、原子のゆれや不純物などの並びの乱れ。}
銅の電子は、原子の約150個分も進んでから散らばります（⑩・例題9）。
\end{tsumazuki}

\begin{tsumazuki}{電子がたくさんあるほど、電気を通す？}
\bkng{電子がぎっしりつまった帯ほど、運び手が多いので、電気をよく通す。}
\bkok{ぎっしりうまった帯では、電子は移る席がなく、電流は流れない。途中までうまった帯が、電気を通す。}
絶縁体の価電子帯にも、電子はたくさんあります（⑧）。
\end{tsumazuki}

\begin{tsumazuki}{温度が上がると、どんな物も抵抗が増える？}
\bkng{温度が上がると、原子がゆれてぶつかりやすくなるので、半導体も抵抗が増える。}
\bkok{半導体では、上の帯へ上がる電子の数が $e^{-E_{\mathrm g}/(2k_{\mathrm B}T)}$ で急に増えるので、抵抗は下がる。}
金属は電子の数がほぼ変わらず、ぶつかりやすさだけが増えるので、抵抗が増えます（⑧の2倍の図）。
\end{tsumazuki}

\begin{tsumazuki}{フェルミエネルギーは、熱のエネルギー？}
\bkng{銅の電子が 7 eV もの大きなエネルギーをもつのは、温度が高いから。}
\bkok{$E_{\mathrm F}$ は、電子が席につまったことで決まり、0 K でも同じ。温度とは関係がない。}
$T_{\mathrm F}=\qty{8.2e4}{K}$ は、$E_{\mathrm F}$ を温度の単位で表しただけで、銅の温度ではありません（④）。
\end{tsumazuki}

%% ======================== 工学ではこう使う ========================
\Needspace{18\baselineskip}\bksection{phgear}{工学ではこう使う}
\begin{phkougaku}[帯とすき間を使い分ける]
\textbf{材料を調べる}：X線の回折（②）は、材料の結晶の並びや、どの物質がまじっているか、金属の中にどれだけの力がかかっているか（原子の面の間かくのわずかなのび）を、こわさずに調べる、工場や研究所の基本の道具です。

\begin{phketa}[固体の電子の数の目安]
\phketaline{格子定数・原子の間かく}{0.2〜0.6 nm}
\phketaline{金属の $E_{\mathrm F}$・$v_{\mathrm F}$}{3〜12 eV、約 $10^6$ m/s（銅 7.0 eV）}
\phketaline{禁制帯幅}{ケイ素 1.12 eV、ガラス 約 9 eV。室温の $k_{\mathrm B}T$ は 0.026 eV}
\phketaline{金属の平均自由行程}{室温で数十 nm（銅 38 nm）}
\phboxhead{答えが桁で合っているか}
固体の電子のエネルギーは、原子の間かく（0.3 nm くらい）の箱の電子と同じくらいで、数 eV です。$E_{\mathrm F}$ が何千 eV や何 meV になったら、$n$ の単位（cm$^{-3}$ と m$^{-3}$ で $10^6$ 倍違う）か、$2/3$ 乗のまちがいを疑いましょう。
\end{phketa}

\textbf{電気を通す・通さない}：電線の銅とアルミニウム（帯が途中まで）、プリント基板や電線の被覆の樹脂やガラス（大きなすき間）、そしてその中間の半導体（ケイ素のウエハー：原-9・原-10）を、目的に合わせて使い分けます。

\textbf{熱を逃がす}：パソコンの CPU の熱を逃がすヒートシンクに銅やアルミニウムを使うのは、電子が熱もよく運ぶからです（ウィーデマン・フランツ：⑩）。

\textbf{超伝導}：MRI の磁石（⑫）、加速器の電磁石、研究用の強い磁石は、Nb-Ti などの超伝導の線を液体ヘリウムで冷やして使います。
\end{phkougaku}
\phstation{固体物理,材料工学,電子材料,超伝導工学}{結晶構造、X線回折、エネルギーバンド、フェルミ準位、禁制帯幅、価電子帯、伝導帯、超伝導}

%% ======================== 章末問題 ========================
\bkendproblems
\bkproblemgroup{基本}{3}
\bkproblem 面の間かく \qty{0.200}{nm}の結晶に、波長 \qty{0.154}{nm}の X線を当てます。$n=1$ と $n=2$ で強く反射する角 $\theta$ を求めましょう。
\bkpans{$n=1$：$\sin\theta=0.385$、$\theta\approx22.6^\circ$。$n=2$：$\sin\theta=0.770$、$\theta\approx50.4^\circ$}
\bkproblem ケイ素の結晶に銅の Kα$_1$ 線（\qty{0.1541}{nm}）を当てると、$2\theta=28.44^\circ$ で強い反射がありました（$n=1$）。この反射をした面の間かくを求めましょう。
\bkpans{$\theta=14.22^\circ$。$d=\lambda/(2\sin\theta)\approx\qty{0.314}{nm}$}
\bkproblem 鉄は体心立方（立方体1つに2個）です。格子定数を \qty{0.287}{nm}、モル質量を \qty{55.85}{g/mol}として、鉄の密度を求めましょう。
\bkpans{$\rho=2M/(N_{\mathrm A}a^3)$ に入れて、約 \qty{7.85}{g/cm^3}}
\bkproblem イオンが間かく \qty{0.20}{nm}で並んだ1次元の結晶で、自由電子の波がはじめてブラッグ反射するときの、電子の波長とエネルギーを求めましょう。
\bkpans{$\lambda=2a=\qty{0.40}{nm}$。$E=h^2/(8m_{\mathrm e}a^2)\approx\qty{9.4}{eV}$}
\bkproblem ナトリウム（原子1個から電子1個）、マグネシウム（2個）、ケイ素（4個で共有結合）のうち、金属になるのはどれですか。帯のうまり方で理由を答えましょう。
\bkpans{ナトリウムは $s$ の帯が半分までうまるので金属。マグネシウムは $s$ の帯がうまりそうだが、$p$ の帯と重なっていて空いた席がすぐ上にあるので金属。ケイ素は価電子帯がぎっしりうまり、1.12 eV のすき間があるので半導体（金属ではない）}
\bkproblemgroup{チャレンジ}{5}
\bkproblem アルミニウムは原子1個から電子を3個出します。密度 \qty{2.70}{g/cm^3}、モル質量 \qty{26.98}{g/mol}として、自由電子の数とフェルミエネルギーを求めましょう。
\bkpans{$n=3\rho N_{\mathrm A}/M\approx1.81\times10^{29}$ 個/m$^3$。$E_{\mathrm F}=\frac{\hbar^2}{2m_{\mathrm e}}(3\pi^2n)^{2/3}\approx\qty{11.7}{eV}$}
\bkproblem 600 K（$k_{\mathrm B}T=\qty{0.0517}{eV}$）で、$E_{\mathrm F}$ より何 eV 上の席で、電子がいる確率が 0.01 になりますか。
\bkpans{$e^{x}+1=100$ より $x=\ln99\approx4.60$。$4.60\times0.0517\approx\qty{0.24}{eV}$ 上}
\bkproblem 銀の抵抗率は 20 ℃で $\qty{1.59e-8}{\ohm.m}$ です（電-4）。ウィーデマン・フランツの法則で、293 K の熱伝導率を求めましょう。
\bkpans{$\kappa=2.44\times10^{-8}\times T/\rho$ で、約 \qty{4.5e2}{W/(m.K)}}
\bkproblem 銅（$T_{\mathrm F}=\qty{8.2e4}{K}$、$\theta_{\mathrm D}=\qty{343}{K}$）で、電子の比熱 $\frac{\pi^2}{2}R\frac{T}{T_{\mathrm F}}$ と、原子のゆれの比熱 $\frac{12\pi^4}{5}R\bigl(\frac{T}{\theta_{\mathrm D}}\bigr)^3$（どちらも 1 mol あたり）が等しくなる温度を求めましょう。
\bkpans{$T^2=\frac{5\,\theta_{\mathrm D}^{\,3}}{24\pi^2T_{\mathrm F}}$ より $T\approx\qty{3.2}{K}$。これより低い温度では電子の比熱のほうが大きい}

\begin{furikaeri}
\begin{checklist}
\item 金属・半導体・絶縁体の違いを、電子のエネルギーの帯（バンド）と、帯のすき間で説明できる。\bkfurref{ph:原-8:sec8}
\item ブラッグの条件で X線の回折の角を求め、格子定数から結晶の密度や原子の数を計算できる。\bkfurref{ph:原-8:sec2}
\item 3次元の箱で電子の席を数えて、金属のフェルミエネルギーとフェルミ速度を求められる。\bkfurref{ph:原-8:sec4}
\item フェルミ分布で、電子が席をうめる割合を求め、金属の電子の比熱が小さいわけを説明できる。\bkfurref{ph:原-8:sec5}
\item くし形の結晶の帯をプログラムで求め、平均自由行程・ウィーデマン・フランツの法則・デバイの $T^3$ 則を使える。\bkfurref{ph:原-8:sec10}
\end{checklist}
\end{furikaeri}
\bkfukidashi{sora}{銅の電子は、0 K でも1秒に1500 km 以上で走り回っているのに、ガラスの電子は、すき間をこえられずに動けないんだね。帯がうまっているかどうかで、こんなに違うなんて！}
